We wish to describe various \enquote{transport} constructions for monoids and modules across different duals of a monoidal category $\cV$. To do this, we first briefly describe these relevant dualities on $\MonCat_{x}$, where $x$ is used as a placeholder for lax, colax, or pseudo. Next, we discuss how morphisms of involutions and inverses provide a natural framework to encode this dualities. For encoding the dualities of modules, virtual double categories shall play an essential role. Moreover, these dualities themselves generalize appropriately to virtual double categories.


\subsection{\tpdf{Dualities of $\MonCat$}} \label{sub:moncatdual}

The $\op_1$ duality (and its generalized inverse $\op_1^t$) can be extended to the $2$-category of monoidal categories and provides the following relation between the lax and colax variants.

\begin{lemma}
	There exists two functors \[\op_1 \colon \op_2(\MonCat_{c}) \rightarrow \MonCat_{l}\] and \[\op_1^t \colon  \op_2(\MonCat_{l}) \rightarrow \MonCat_{c}\] extending the involution $\op_1 \colon \op_2(\nCat{1}) \rightarrow \nCat{1}$ along the forgetful $2$-functors $U_{x} \colon\MonCat_{x} \rightarrow \Cat$. These functors are strict inverses of each other relative to the $\op_2$ involution. That is, they satisfy the constraints as defined by the following diagrams:
	% https://q.uiver.app/#q=WzAsNixbMCwwLCJcXE1vbkNhdF9jIl0sWzEsMSwiXFxvcF8yKFxcTW9uQ2F0X2wpIl0sWzAsMiwiXFxNb25DYXRfYyJdLFsyLDAsIlxcTW9uQ2F0X2wiXSxbMywxLCJcXG9wXzIoXFxNb25DYXRfYykiXSxbMiwyLCJcXE1vbkNhdF9sIl0sWzAsMSwiXFxvcF8yKFxcb3BfMSkiXSxbMSwyLCJcXG9wXnRfMSJdLFswLDIsIjEiLDJdLFszLDQsIlxcb3BfMihcXG9wXnRfMSkiXSxbNCw1LCJcXG9wXzEiXSxbMyw1LCIxIiwyXV0=
	\[\begin{tikzcd}[ampersand replacement=\&]
		{\MonCat_c} \&\& {\MonCat_l} \& \\
		\& {\op_2(\MonCat_l)} \&\& {\op_2(\MonCat_c)} \\
		{\MonCat_c} \&\& {\MonCat_l}
		\arrow["{\op_2(\op_1)}", from=1-1, to=2-2]
		\arrow["1"', from=1-1, to=3-1]
		\arrow["{\op_2(\op^t_1)}", from=1-3, to=2-4]
		\arrow["1"', from=1-3, to=3-3]
		\arrow["{\op^t_1}", from=2-2, to=3-1]
		\arrow["{\op_1}", from=2-4, to=3-3]
	\end{tikzcd}\cdpunct[4pt]{.}\]
	Furthermore, $U_x$ can be extended to a (strict) morphism of inverses and the $2$-functors $\op_1$ and $\op_1^t$ are isomorphic when restricted onto the variant $\MonCat_{p}$.
\end{lemma}

\begin{proof}
	For each diagram defining some structural monoidal data or coherence constraint, apply the functor $\op_1 \colon \op_2(\nCat{1}) \rightarrow \nCat{1}$ to obtain the required dual diagram. More precisely, $\op_1$ is a monoidal $\nCat{1}$-enriched functor of monoidal $\nCat{1}$-enriched categories. Thus, one could also adapt \cite[Thm.~4.5.1]{yauinvcat2020} in our situation by replacing the strict structures with pseudo structures.
\end{proof}

The same duality holds for $\BrMonCat_x$ and $\SymMonCat_x$. For instance, the restriction of $\op_1$ functor yields the commuting diagram
% https://q.uiver.app/#q=WzAsNixbNCwwLCJcXG9wXzIoXFxNb25DYXRfYykiXSxbNCwyLCJcXE1vbkNhdF9sIl0sWzIsMCwiXFxvcF8yKFxcQnJNb25DYXRfYykiXSxbMiwyLCJcXEJyTW9uQ2F0X2wiXSxbMCwwLCJcXG9wXzIoXFxTeW1Nb25DYXRfYykiXSxbMCwyLCJcXFN5bU1vbkNhdF9sIl0sWzAsMSwiXFxvcF8xIl0sWzMsMSwiVSIsMl0sWzIsMCwiXFxvcF8yKFUpIl0sWzIsMywiXFxvcF8xIiwyXSxbNCwyLCJcXG9wXzIoVSkiXSxbNSwzLCJVIiwyXSxbNCw1LCJcXG9wXzEiLDJdXQ==
\[\begin{tikzcd}[ampersand replacement=\&]
	{\op_2(\SymMonCat_c)} \&\& {\op_2(\BrMonCat_c)} \&\& {\op_2(\MonCat_c)} \\
	\\
	{\SymMonCat_l} \&\& {\BrMonCat_l} \&\& {\MonCat_l}
	\arrow["{\op_2(U)}", from=1-1, to=1-3]
	\arrow["{\op_1}"', from=1-1, to=3-1]
	\arrow["{\op_2(U)}", from=1-3, to=1-5]
	\arrow["{\op_1}"', from=1-3, to=3-3]
	\arrow["{\op_1}", from=1-5, to=3-5]
	\arrow["U"', from=3-1, to=3-3]
	\arrow["U"', from=3-3, to=3-5]
\end{tikzcd}\cdpunct[8pt]{.}\]
Here, $\op_1$ maps a braided monoidal category as $(\cV, b_{x,y}) \mapsto (\op_1(\cV),b^{-1}_{x,y})$. And, each square in the diagram represents a part of the corresponding morphism of inverses.

\begin{remark}
	We can combine both $\op_1$ and $\op_1^t$ as a single involution on the double category of monoidal category (\cite[Sec.~4.2.6]{grandis2019dblcat}), when considered relative to the transpose involution (\cite[Sec.~3.2.2]{grandis2019dblcat}) of the \emph{$1$-category} of (strict) double categories. In particular, there exists a double functor $\op_1 \colon t(\mathsf{\mathbb{M}onCat}) \rightarrow \mathsf{\mathbb{M}onCat}$ that satisfies $t(\op_1) \circ \op_1 = 1$ or equivalently $\op_1 \circ t(\op_1) = 1$. In particular, we have the following equivalent diagram of morphisms of involutions:
	% https://q.uiver.app/#q=WzAsNixbNCwwLCJ0KFxcbWF0aHNme1xcbWF0aGJie019b25DYXR9KSJdLFs0LDIsIlxcbWF0aHNme1xcbWF0aGJie019b25DYXR9Il0sWzIsMCwidChcXG1hdGhzZntCclxcbWF0aGJie019b25DYXR9KSJdLFsyLDIsIlxcbWF0aHNme0JyXFxtYXRoYmJ7TX1vbkNhdH0iXSxbMCwwLCJ0KFxcbWF0aHNme1N5bVxcbWF0aGJie019b25DYXR9KSJdLFswLDIsIlxcbWF0aHNme1N5bVxcbWF0aGJie019b25DYXR9Il0sWzAsMSwiXFxvcF8xIl0sWzMsMSwiVSIsMl0sWzIsMCwidChVKSJdLFsyLDMsIlxcb3BfMSIsMl0sWzQsMiwidChVKSJdLFs1LDMsIlUiLDJdLFs0LDUsIlxcb3BfMSIsMl1d
	\[\begin{tikzcd}[ampersand replacement=\&]
		{t(\mathsf{Sym\mathbb{M}onCat})} \&\& {t(\mathsf{Br\mathbb{M}onCat})} \&\& {t(\mathsf{\mathbb{M}onCat})} \\
		\\
		{\mathsf{Sym\mathbb{M}onCat}} \&\& {\mathsf{Br\mathbb{M}onCat}} \&\& {\mathsf{\mathbb{M}onCat}}
		\arrow["{t(U)}", from=1-1, to=1-3]
		\arrow["{\op_1}"', from=1-1, to=3-1]
		\arrow["{t(U)}", from=1-3, to=1-5]
		\arrow["{\op_1}"', from=1-3, to=3-3]
		\arrow["{\op_1}", from=1-5, to=3-5]
		\arrow["U"', from=3-1, to=3-3]
		\arrow["U"', from=3-3, to=3-5]
	\end{tikzcd}\cdpunct[8pt]{.}\]
\end{remark}


The other important duality is that of reversal of a monoidal category. In particular, for each monoidal category $(\cV, \otimes ,I)$, there exists a corresponding \emph{reversed monoidal category} $(\cV, \otimes_r,I)$, which is denoted by $r(\cV)$. The monoidal product of $r(\cV)$ is defined using the composite $\cV \times \cV \xrightarrow{s} \cV \times \cV \xrightarrow{\otimes} \cV$, and the monoidal unit remains the same. The associator is given by $\left(a_{r(\cV)}\right)_{x,y,z}= \left(a_{\cV}\right)^{-1}_{z,y,x}$, and the unitors by $\lambda_{r(\cV)} = \rho_{\cV}$, $\rho_{r(\cV)} = \lambda_{\cV}$.

The reversal construction $r$ also extends to an involution. That is, $r \colon \MonCat_{x} \rightarrow \MonCat_{x}$ is a $2$-functor which satisfy $r^2=1$. Furthermore, this duality restricts to the $2$-categories $\BrMonCat_x$ and $\SymMonCat_x$ appropriately, as before. That is, there are commuting squares

% https://q.uiver.app/#q=WzAsNixbNCwwLCJcXE1vbkNhdF94Il0sWzQsMiwiXFxNb25DYXRfeCJdLFsyLDAsIlxcQnJNb25DYXRfeCJdLFsyLDIsIlxcQnJNb25DYXRfeCJdLFswLDAsIlxcU3ltTW9uQ2F0X3giXSxbMCwyLCJcXFN5bU1vbkNhdF94Il0sWzAsMSwiciJdLFszLDEsIlUiLDJdLFsyLDAsIlUiXSxbMiwzLCJyIiwyXSxbNCwyLCJVIl0sWzUsMywiVSIsMl0sWzQsNSwiciIsMl1d
\[\begin{tikzcd}[ampersand replacement=\&]
	{\SymMonCat_x} \&\& {\BrMonCat_x} \&\& {\MonCat_x} \\
	\\
	{\SymMonCat_x} \&\& {\BrMonCat_x} \&\& {\MonCat_x}
	\arrow["U", from=1-1, to=1-3]
	\arrow["r"', from=1-1, to=3-1]
	\arrow["U", from=1-3, to=1-5]
	\arrow["r"', from=1-3, to=3-3]
	\arrow["r", from=1-5, to=3-5]
	\arrow["U"', from=3-1, to=3-3]
	\arrow["U"', from=3-3, to=3-5]
\end{tikzcd}\cdpunct[8pt]{,}\]
representing the corresponding morphisms of involutions. Here, $r$ sends a braiding $b_{x,y}$ to the reversed braiding $b_{y,x}$.

\begin{remark}
	Analogous double categorical extension to $\mathsf{\mathbb{M}onCat}$ (and variants) can also be formulated for the reversal duality as relative to the identity involution of strict double categories. 
\end{remark}

\begin{comment}
	\subsubsection{\tpdf{$i$ duality involution}}
	
	If $\cV$ is braided with braiding $b_{x,y}$, we have one more nontrivial duality that is not detectable by the forgetful functor $U \colon \BrMonCat_{x} \rightarrow \MonCat_{x}$. This can be  observed from the following four possible choices of braided monoidal structures, as seen below, all connected by the identity functor $1_{\cV}$:
	% https://q.uiver.app/#q=WzAsNSxbNCwwLCIoXFxjVixiX3t4LHl9KSJdLFswLDAsIihyKFxcY1YpLGJfe3kseH0pIl0sWzQsMiwiKFxcY1YsYl57LTF9X3t5LHh9KSJdLFswLDIsIihyKFxcY1YpLGJeey0xfV97eCx5fSkiXSxbMiwxLCJcXHRleHR7YnJhaWRlZCBpZiAkXFxjViQgaXMgc3ltbWV0cmljfSJdLFswLDEsIlxcdGV4dHticmFpZGVkIG1vbn0iLDIseyJzdHlsZSI6eyJoZWFkIjp7Im5hbWUiOiJub25lIn19fV0sWzAsMiwiXFx0ZXh0e21vbn0iLDAseyJzdHlsZSI6eyJoZWFkIjp7Im5hbWUiOiJub25lIn19fV0sWzEsMywiXFx0ZXh0e21vbn0iLDIseyJzdHlsZSI6eyJoZWFkIjp7Im5hbWUiOiJub25lIn19fV0sWzIsMywiXFx0ZXh0e2JyYWlkZWQgbW9ufSIsMCx7InN0eWxlIjp7ImhlYWQiOnsibmFtZSI6Im5vbmUifX19XSxbNCw3LCIiLDIseyJzaG9ydGVuIjp7InRhcmdldCI6MjB9LCJsZXZlbCI6MSwic3R5bGUiOnsiYm9keSI6eyJuYW1lIjoic3F1aWdnbHkifX19XSxbNCw2LCIiLDAseyJzaG9ydGVuIjp7InRhcmdldCI6MjB9LCJsZXZlbCI6MSwic3R5bGUiOnsiYm9keSI6eyJuYW1lIjoic3F1aWdnbHkifX19XV0=
	\[\begin{tikzcd}[ampersand replacement=\&]
		{(r(\cV),b_{y,x})} \&\&\&\& {(\cV,b_{x,y})} \\
		\&\& {\text{braided if $\cV$ is symmetric}} \\
		{(r(\cV),b^{-1}_{x,y})} \&\&\&\& {(\cV,b^{-1}_{y,x})}
		\arrow[""{name=0, anchor=center, inner sep=0}, "{\text{mon}}"', no head, from=1-1, to=3-1]
		\arrow["{\text{braided mon}}"', no head, from=1-5, to=1-1]
		\arrow[""{name=1, anchor=center, inner sep=0}, "{\text{mon}}", no head, from=1-5, to=3-5]
		\arrow["{\text{braided mon}}", no head, from=3-5, to=3-1]
		\arrow[between={0}{0.8}, squiggly, from=2-3, to=0]
		\arrow[between={0}{0.8}, squiggly, from=2-3, to=1]
	\end{tikzcd}\cdpunct[8pt]{.}\]
	
	In particular, we obtain one more duality $i \colon \BrMonCat_x \rightarrow \BrMonCat_x$, called the inversion duality, given by $b_{x,y} \mapsto b^{-1}_{y,x}$. On symmetric monoidal categories, this evaluates to just the identity functor. In particular, the following analogous diagram commute:
	% https://q.uiver.app/#q=WzAsNixbMCwwLCJcXFN5bU1vbkNhdF94Il0sWzAsMiwiXFxTeW1Nb25DYXRfeCJdLFs0LDAsIlxcQnJNb25DYXRfeCJdLFs0LDIsIlxcQnJNb25DYXRfeCJdLFs4LDAsIlxcTW9uQ2F0X3giXSxbOCwyLCJcXE1vbkNhdF94Il0sWzIsMywiaSIsMl0sWzAsMSwiMSIsMl0sWzAsMiwiVSJdLFsxLDMsIlUiLDJdLFs0LDUsIjEiXSxbMiw0LCJVIl0sWzMsNSwiVSIsMl1d
	\[\begin{tikzcd}[ampersand replacement=\&]
		{\SymMonCat_x} \&\&\&\& {\BrMonCat_x} \&\&\&\& {\MonCat_x} \\
		\\
		{\SymMonCat_x} \&\&\&\& {\BrMonCat_x} \&\&\&\& {\MonCat_x}
		\arrow["U", from=1-1, to=1-5]
		\arrow["1"', from=1-1, to=3-1]
		\arrow["U", from=1-5, to=1-9]
		\arrow["i"', from=1-5, to=3-5]
		\arrow["1", from=1-9, to=3-9]
		\arrow["U"', from=3-1, to=3-5]
		\arrow["U"', from=3-5, to=3-9]
	\end{tikzcd}\cdpunct[8pt]{.}\]
	
	\begin{remark}
		As expected, there are also double categorical extension to $\mathsf{\mathbb{M}onCat}$ (and variants) for the inversion duality. 
	\end{remark}
	
\end{comment}

\subsection{Transport of monoids in virtual double categories} \label{sub:monmoddualvdc}

It is well known that the category of monoids in a monoidal category can be identified using the isomorphisms $\Mon(\cV) \cong \MonCat_l(*,\cV)$. Similarly, commutative monoids in a braided monoidal category can be identified using $\cMon(\cV)\cong \BrMonCat_l(*,\cV)$. And if $\cV$ is symmetric, we have further isomorphisms $\BrMonCat_l(*,\cV) \cong \SymMonCat_l(*,\cV)$. Our goal is to understand various ways to transport monoids and comonoids using these identifications. As mentioned before, we also wish to describe these constructions within the framework of generalized involutions/inverses, with the main motivation being the intuition for what to expect in higher dimensions.  

We start with defining the $\op$ duality functors between monoids and comonoids for a given $\cV$. 

\begin{definition}
	For each monoidal category $\cV$, define a $1$-functor $\op_{\cV} \colon \op_1(\CoMon(\cV)) \rightarrow \Mon(\op_1(\cV))$ using the commutative diagram
	% https://q.uiver.app/#q=WzAsNCxbMiwxLCJcXE1vbkNhdF9sKCosXFxvcF8xKFxcY1YpKSJdLFsyLDAsIlxcTW9uKFxcb3BfMShcXGNWKSkiXSxbMCwwLCJcXG9wXzEoXFxDb01vbihcXGNWKSkiXSxbMCwxLCJcXG9wXzEoXFxNb25DYXRfYygqLFxcY1YpKSJdLFsyLDEsIlxcb3Bfe1xcY1Z9Il0sWzMsMCwiXFxvcF8xIChcXGNvbmcpIiwyXSxbMiwzLCJcXGNvbmciLDAseyJzdHlsZSI6eyJoZWFkIjp7Im5hbWUiOiJub25lIn19fV0sWzEsMCwiXFxjb25nIiwyLHsic3R5bGUiOnsiaGVhZCI6eyJuYW1lIjoibm9uZSJ9fX1dXQ==
	\[\begin{tikzcd}[ampersand replacement=\&]
		{\op_1(\CoMon(\cV))} \&\& {\Mon(\op_1(\cV))} \\
		{\op_1(\MonCat_c(*,\cV))} \&\& {\MonCat_l(*,\op_1(\cV))}
		\arrow["{\op_{\cV}}", from=1-1, to=1-3]
		\arrow["\cong", no head, from=1-1, to=2-1]
		\arrow["\cong"', no head, from=1-3, to=2-3]
		\arrow["{\op_1 (\cong)}"', from=2-1, to=2-3]
	\end{tikzcd}\cdpunct[8pt]{.}\]
	The dual version $\op^t_{\cV} \colon \op_1(\Mon(\cV)) \rightarrow \CoMon(\op_1(\cV))$ is defined similarly.
\end{definition}

We already observe that there is no obvious way to compose $\op_{\cV}$ and $\op^t_{\cV}$ to produce the identity map as all inverses do. At the same time, it is clear that this is a form of duality. The primary issue is that we are not incorporating the functoriality in $\cV$. In particular, both $\CoMon(-)$ and $\Mon(-)$ are strict 2-functors as they are defined by a corepresentation. Therefore, we ought to extend both $\op_{\cV}$ and $\op^{t}_{\cV}$ to strict $2$-natural transformations.

\begin{lemma}
	The functors $\op_{\cV}$ and $\op^t_{\cV}$ form the components of strict $2$-natural transformations
	% https://q.uiver.app/#q=WzAsOCxbMCwwLCJcXG9wXzIoXFxNb25DYXRfYykiXSxbMywwLCJcXG9wXzIoXFxDYXQpIl0sWzMsMywiXFxDYXQiXSxbMCwzLCJcXE1vbkNhdF9sIl0sWzQsMCwiXFxvcF8yKFxcTW9uQ2F0X2wpIl0sWzQsMywiXFxNb25DYXRfYyJdLFs3LDAsIlxcb3BfMihcXENhdCkiXSxbNywzLCJcXENhdCJdLFswLDEsIlxcb3BfMihcXENvTW9uKC0pKSJdLFsxLDIsIlxcb3BfMSJdLFswLDMsIlxcb3BfMSIsMl0sWzMsMiwiXFxNb24oLSkiLDJdLFsxLDMsIlxcb3AiLDIseyJzaG9ydGVuIjp7InNvdXJjZSI6MjAsInRhcmdldCI6MjB9LCJsZXZlbCI6Mn1dLFs0LDUsIlxcb3BedF8xIiwyXSxbNiw3LCJcXG9wXzEiXSxbNCw2LCJcXG9wXzIoXFxNb24oLSkpIl0sWzUsNywiXFxDb01vbigtKSIsMl0sWzYsNSwiXFxvcF50IiwyLHsic2hvcnRlbiI6eyJzb3VyY2UiOjIwLCJ0YXJnZXQiOjIwfSwibGV2ZWwiOjJ9XV0=
	\[\begin{tikzcd}[ampersand replacement=\&,column sep=scriptsize]
		{\op_2(\MonCat_c)} \&\&\& {\op_2(\Cat)} \& {\op_2(\MonCat_l)} \&\&\& {\op_2(\Cat)} \\
		\\
		\\
		{\MonCat_l} \&\&\& \Cat \& {\MonCat_c} \&\&\& \Cat
		\arrow["{\op_2(\CoMon(-))}", from=1-1, to=1-4]
		\arrow["{\op_1}"', from=1-1, to=4-1]
		\arrow["\op"', between={0.2}{0.8}, Rightarrow, nfold, from=1-4, to=4-1]
		\arrow["{\op_1}", from=1-4, to=4-4]
		\arrow["{\op_2(\Mon(-))}", from=1-5, to=1-8]
		\arrow["{\op^t_1}"', from=1-5, to=4-5]
		\arrow["{\op^t}"', between={0.2}{0.8}, Rightarrow, nfold, from=1-8, to=4-5]
		\arrow["{\op_1}", from=1-8, to=4-8]
		\arrow["{\Mon(-)}"', from=4-1, to=4-4]
		\arrow["{\CoMon(-)}"', from=4-5, to=4-8]
	\end{tikzcd}\cdpunct[8pt]{.}\]
	Furthermore, these two transformations form a morphism of inverses relative to the strict $2$-involution $\op_2 \colon \nCat{2}_{\leq 2} \rightarrow \nCat{2}_{\leq 2}$.
\end{lemma}

\begin{proof}
	The $2$-naturality and strictness of the two collection of maps $\op_{\cV}$ and $\op_{\cV}^{t}$ is a direct consequence of their definition in terms of the corepresentation by the terminal monoidal category. 
	
	The two axioms for morphism of inverses are similarly straightforward. For instance, one of these axioms can be verified as indicated below:
	% https://q.uiver.app/#q=WzAsMTUsWzAsMiwiXFxvcF8yXjIoXFxNb25DYXRfYykiXSxbMiwyLCJcXG9wXzJeMihcXENhdCkiXSxbMCw0LCJcXG9wXzIoXFxNb25DYXRfbCkiXSxbMSw2LCJcXE1vbkNhdF9jIl0sWzIsNCwiXFxvcF8yKFxcQ2F0KSJdLFszLDYsIlxcQ2F0Il0sWzAsOSwiXFxDb01vbihcXGNWKSJdLFswLDcsIlxcb3BfMShcXE1vbihcXG9wXzEoXFxjVikpKSJdLFs1LDcsIlxcQ29Nb24oXFxvcF8xXjIoXFxjVikpIl0sWzUsOSwiXFxDb01vbihcXGNWKSJdLFszLDAsIlxcQ2F0Il0sWzEsMCwiXFxNb25DYXRfYyJdLFs1LDAsIlxcTW9uQ2F0X2MiXSxbNSw2LCJcXENhdCJdLFs0LDMsIlxcZXF1aXYiXSxbMCwxLCJcXG9wXzJeMihcXENvTW9uKC0pKSJdLFsyLDMsIlxcb3BedF8xIiwyXSxbNCw1LCJcXG9wXzEiLDFdLFsyLDQsIlxcb3BfMihcXE1vbigtKSkiXSxbMyw1LCJcXENvTW9uKC0pIiwyXSxbNCwzLCJcXG9wXnQiLDIseyJzaG9ydGVuIjp7InNvdXJjZSI6MjAsInRhcmdldCI6MjB9LCJsZXZlbCI6Mn1dLFsxLDQsIlxcb3BfMihcXG9wXzEpIiwxXSxbMCwyLCJcXG9wXzIoXFxvcF8xKSIsMl0sWzEsMiwiXFxvcF8yKFxcb3ApIiwyLHsic2hvcnRlbiI6eyJzb3VyY2UiOjIwLCJ0YXJnZXQiOjIwfSwibGV2ZWwiOjJ9XSxbNiw3LCJcXG9wX3tcXGNWfSJdLFs3LDgsIlxcb3BedF97XFxvcF8xe1xcY1Z9fSJdLFs4LDksIlxcQ29Nb24oMV97XFxjVn0pIl0sWzYsOSwiMV97XFxDb01vbihcXGNWKX0iLDJdLFsxMCw1LCIxX3tcXENhdH0iXSxbMTAsMSwiIiwwLHsic3R5bGUiOnsiaGVhZCI6eyJuYW1lIjoibm9uZSJ9fX1dLFsxMSwwLCIiLDAseyJzdHlsZSI6eyJoZWFkIjp7Im5hbWUiOiJub25lIn19fV0sWzExLDEwLCJcXENvTW9uKC0pIl0sWzEwLDAsIiIsMCx7InNob3J0ZW4iOnsic291cmNlIjoyMCwidGFyZ2V0IjoyMH0sImxldmVsIjoyLCJzdHlsZSI6eyJoZWFkIjp7Im5hbWUiOiJub25lIn19fV0sWzEyLDEzLCJcXENvTW9uKC0pIiwxXSxbMjEsMjgsIiIsMSx7InNob3J0ZW4iOnsic291cmNlIjozMCwidGFyZ2V0IjoyMH0sInN0eWxlIjp7ImhlYWQiOnsibmFtZSI6Im5vbmUifX19XV0=
	\[\begin{tikzcd}[ampersand replacement=\&,cramped,column sep=small]
		\& {\MonCat_c} \&\& \Cat \&\& {\MonCat_c} \\
		\\
		{\op_2^2(\MonCat_c)} \&\& {\op_2^2(\Cat)} \\
		\&\&\&\& \equiv \\
		{\op_2(\MonCat_l)} \&\& {\op_2(\Cat)} \\
		\\
		\& {\MonCat_c} \&\& \Cat \&\& \Cat \\
		{\op_1(\Mon(\op_1(\cV)))} \&\&\&\&\& {\CoMon(\op_1^2(\cV))} \\
		\\
		{\CoMon(\cV)} \&\&\&\&\& {\CoMon(\cV)}
		\arrow["{\CoMon(-)}", from=1-2, to=1-4]
		\arrow[no head, from=1-2, to=3-1]
		\arrow[between={0.2}{0.8}, equals, nfold, from=1-4, to=3-1]
		\arrow[no head, from=1-4, to=3-3]
		\arrow[""{name=0, anchor=center, inner sep=0}, "{1_{\Cat}}", from=1-4, to=7-4]
		\arrow["{\CoMon(-)}"{description}, from=1-6, to=7-6]
		\arrow["{\op_2^2(\CoMon(-))}", from=3-1, to=3-3]
		\arrow["{\op_2(\op_1)}"', from=3-1, to=5-1]
		\arrow["{\op_2(\op)}"', between={0.2}{0.8}, Rightarrow, nfold, from=3-3, to=5-1]
		\arrow[""{name=1, anchor=center, inner sep=0}, "{\op_2(\op_1)}"{description}, from=3-3, to=5-3]
		\arrow["{\op_2(\Mon(-))}", from=5-1, to=5-3]
		\arrow["{\op^t_1}"', from=5-1, to=7-2]
		\arrow["{\op^t}"', between={0.2}{0.8}, Rightarrow, nfold, from=5-3, to=7-2]
		\arrow["{\op_1}"{description}, from=5-3, to=7-4]
		\arrow["{\CoMon(-)}"', from=7-2, to=7-4]
		\arrow["{\op^t_{\op_1{\cV}}}", from=8-1, to=8-6]
		\arrow["{\CoMon(1_{\cV})}", from=8-6, to=10-6]
		\arrow["{\op_{\cV}}", from=10-1, to=8-1]
		\arrow["{1_{\CoMon(\cV)}}"', from=10-1, to=10-6]
		\arrow[between={0.3}{0.8}, equals, nfold, from=1, to=0]
	\end{tikzcd}\cdpunct[8pt]{.}\]
\end{proof}

\begin{comment}
	If $\cV$ is either a braided or a symmetric monoidal category, we can similarly extend these duality natural transformations to commutative monoids and comonoids. In fact, the commutative variant of the dualities have canonical embeddings into the general one through strict $2$-transformations
	% https://q.uiver.app/#q=WzAsNCxbNiwwLCJcXE1vbkNhdF9sIl0sWzMsMCwiXFxCck1vbkNhdF9sIl0sWzAsMCwiXFxTeW1Nb25DYXRfbCJdLFszLDQsIlxcQ2F0Il0sWzEsMCwiVSJdLFsyLDEsIlUiXSxbMCwzLCJcXE1vbigtKSIsMV0sWzEsMywiXFxjTW9uKC0pIiwxXSxbMiwzLCJcXGNNb24oLSkiLDFdLFs3LDYsIiIsMSx7InNob3J0ZW4iOnsic291cmNlIjozMCwidGFyZ2V0IjozMH19XSxbOCw3LCIiLDEseyJzaG9ydGVuIjp7InNvdXJjZSI6MzAsInRhcmdldCI6MzB9LCJzdHlsZSI6eyJoZWFkIjp7Im5hbWUiOiJub25lIn19fV1d
	\[\begin{tikzcd}[ampersand replacement=\&,cramped]
		{\SymMonCat_l} \&\&\& {\BrMonCat_l} \&\&\& {\MonCat_l} \\
		\\
		\\
		\\
		\&\&\& \Cat
		\arrow["U", from=1-1, to=1-4]
		\arrow[""{name=0, anchor=center, inner sep=0}, "{\cMon(-)}"{description}, from=1-1, to=5-4]
		\arrow["U", from=1-4, to=1-7]
		\arrow[""{name=1, anchor=center, inner sep=0}, "{\cMon(-)}"{description}, from=1-4, to=5-4]
		\arrow[""{name=2, anchor=center, inner sep=0}, "{\Mon(-)}"{description}, from=1-7, to=5-4]
		\arrow[between={0.3}{0.7}, equals, nfold, from=0, to=1]
		\arrow[between={0.3}{0.7}, Rightarrow, nfold, from=1, to=2]
	\end{tikzcd}\]
	
	% https://q.uiver.app/#q=WzAsNCxbMCwwLCJcXFN5bU1vbkNhdF9jIl0sWzMsMCwiXFxCck1vbkNhdF9jIl0sWzYsMCwiXFxNb25DYXRfYyJdLFszLDQsIlxcQ2F0Il0sWzIsMywiXFxDb01vbigtKSIsMV0sWzEsMiwiVSJdLFswLDEsIlUiXSxbMSwzLCJcXGNDb01vbigtKSIsMV0sWzAsMywiXFxjQ29Nb24oLSkiLDFdLFs4LDcsIiIsMSx7InNob3J0ZW4iOnsic291cmNlIjozMCwidGFyZ2V0IjozMH0sInN0eWxlIjp7ImhlYWQiOnsibmFtZSI6Im5vbmUifX19XSxbNyw0LCIiLDEseyJzaG9ydGVuIjp7InNvdXJjZSI6MzAsInRhcmdldCI6MzB9fV1d
	\[\begin{tikzcd}[ampersand replacement=\&,cramped]
		{\SymMonCat_c} \&\&\& {\BrMonCat_c} \&\&\& {\MonCat_c} \\
		\\
		\\
		\\
		\&\&\& \Cat
		\arrow["U", from=1-1, to=1-4]
		\arrow[""{name=0, anchor=center, inner sep=0}, "{\cCoMon(-)}"{description}, from=1-1, to=5-4]
		\arrow["U", from=1-4, to=1-7]
		\arrow[""{name=1, anchor=center, inner sep=0}, "{\cCoMon(-)}"{description}, from=1-4, to=5-4]
		\arrow[""{name=2, anchor=center, inner sep=0}, "{\CoMon(-)}"{description}, from=1-7, to=5-4]
		\arrow[between={0.3}{0.7}, equals, nfold, from=0, to=1]
		\arrow[between={0.3}{0.7}, Rightarrow, nfold, from=1, to=2]
	\end{tikzcd}\]
	Furthermore, these two embeddings commute appropriately with the $2$-transformations $\op$ and $\op^t$ of the $\op$ duality.
\end{comment}

The phenomenon of transport across $r$ and $i$ is more interesting as considering the braiding or symmetry provides additional transport structure unlike $\op_1$. First, consider that $\cV$ is just an arbitrary monoidal category. In this case, we only have the expected reversal duality.

\begin{lemma}
	There exists a strict 2-natural transformation
	% https://q.uiver.app/#q=WzAsNCxbMCwwLCJcXE1vbkNhdF9sIl0sWzIsMCwiXFxDYXQiXSxbMiwyLCJcXENhdCJdLFswLDIsIlxcTW9uQ2F0X2wiXSxbMCwxLCJcXE1vbigtKSJdLFsxLDIsIjEiXSxbMCwzLCJyIiwyXSxbMywyLCJcXE1vbigtKSIsMl0sWzEsMywiciIsMix7InNob3J0ZW4iOnsic291cmNlIjoyMCwidGFyZ2V0IjoyMH0sImxldmVsIjoyfV1d
	\[\begin{tikzcd}[ampersand replacement=\&]
		{\MonCat_l} \&\& \Cat \\
		\\
		{\MonCat_l} \&\& \Cat
		\arrow["{\Mon(-)}", from=1-1, to=1-3]
		\arrow["r"', from=1-1, to=3-1]
		\arrow["r"', between={0.2}{0.8}, Rightarrow, from=1-3, to=3-1]
		\arrow["1", from=1-3, to=3-3]
		\arrow["{\Mon(-)}"', from=3-1, to=3-3]
	\end{tikzcd}\cdpunct[8pt]{,}\]
	which defines a morphism of involutions relative to the identity functor $1 \colon \nCat{2}_{\leq 2} \rightarrow \nCat{2}_{\leq 2}$.
\end{lemma}

Now take $\cV$ to be a braided monoidal category. The key distinction from $\op_1$ duality is that $\BrMonCat$ admits two distinct non-trivial $2$-cells 
% https://q.uiver.app/#q=WzAsOCxbMCwwLCJcXEJyTW9uQ2F0X2wiXSxbMiwwLCJcXEJyTW9uQ2F0X2wiXSxbMCwxLCJiX3tcXGNWfSA6KHIoXFxjViksYl97eSx4fSkiXSxbMiwxLCIoXFxjVixiX3t4LHl9KSJdLFszLDAsIlxcQnJNb25DYXRfbCJdLFs1LDAsIlxcQnJNb25DYXRfbCJdLFszLDEsImJeey0xfV97XFxjVn0gOihyKFxcY1YpLGJfe3kseH0pIl0sWzUsMSwiKFxcY1YsYl97eCx5fSkiXSxbMCwxLCJyIiwyLHsib2Zmc2V0IjozLCJjdXJ2ZSI6MX1dLFswLDEsIjEiLDAseyJvZmZzZXQiOi0zLCJjdXJ2ZSI6LTF9XSxbMiwzLCIoMSxiX3t4LHl9LDEpIiwyXSxbNCw1LCIxIiwwLHsib2Zmc2V0IjotMywiY3VydmUiOi0xfV0sWzQsNSwiciIsMix7Im9mZnNldCI6MywiY3VydmUiOjF9XSxbNiw3LCIoMSxiXnstMX1fe3kseH0sMSkiLDJdLFs4LDksImIiLDIseyJzaG9ydGVuIjp7InNvdXJjZSI6MjAsInRhcmdldCI6MjB9fV0sWzEyLDExLCJiXnstMX0iLDIseyJzaG9ydGVuIjp7InNvdXJjZSI6MjAsInRhcmdldCI6MjB9fV1d
\[\begin{tikzcd}[ampersand replacement=\&,cramped]
	{\BrMonCat_l} \&\& {\BrMonCat_l} \& {\BrMonCat_l} \&\& {\BrMonCat_l} \\
	{b_{\cV} :(r(\cV),b_{y,x})} \&\& {(\cV,b_{x,y})} \& {b^{-1}_{\cV} :(r(\cV),b_{y,x})} \&\& {(\cV,b_{x,y})}
	\arrow[""{name=0, anchor=center, inner sep=0}, "r"', shift right=3, curve={height=6pt}, from=1-1, to=1-3]
	\arrow[""{name=1, anchor=center, inner sep=0}, "1", shift left=3, curve={height=-6pt}, from=1-1, to=1-3]
	\arrow[""{name=2, anchor=center, inner sep=0}, "1", shift left=3, curve={height=-6pt}, from=1-4, to=1-6]
	\arrow[""{name=3, anchor=center, inner sep=0}, "r"', shift right=3, curve={height=6pt}, from=1-4, to=1-6]
	\arrow["{(1,b_{x,y},1)}"', from=2-1, to=2-3]
	\arrow["{(1,b^{-1}_{y,x},1)}"', from=2-4, to=2-6]
	\arrow["b"', between={0.2}{0.8}, Rightarrow, nfold, from=0, to=1]
	\arrow["{b^{-1}}"', between={0.2}{0.8}, Rightarrow, nfold, from=3, to=2]
\end{tikzcd}\]
between the identity functor and the reversal functor. And, these $2$-cell $b$ and $b^{-1}$ are part of a morphism of inverses, albeit degenerately, 
% https://q.uiver.app/#q=WzAsNixbMywwLCJcXEJyTW9uQ2F0X2wiXSxbMywyLCJcXEJyTW9uQ2F0X2wiXSxbMSwwLCIoXFxjVixiX3t4LHl9KSJdLFswLDIsIihyKFxcY1YpLGJfe3kseH0pIl0sWzMsNCwiXFxCck1vbkNhdF97bH0iXSxbMSw0LCIoXFxjVixiX3t4LHl9KSJdLFswLDEsInIiLDAseyJvZmZzZXQiOi0zLCJjdXJ2ZSI6LTF9XSxbMCwxLCIxIiwyLHsib2Zmc2V0IjozLCJjdXJ2ZSI6MX1dLFsyLDMsInIoYl97XFxjVn0pXFxjb2xvbigxLGJfe3kseH0sMSkiLDJdLFsxLDQsIjEiLDIseyJvZmZzZXQiOjMsImN1cnZlIjoxfV0sWzEsNCwiciIsMCx7Im9mZnNldCI6LTMsImN1cnZlIjotMX1dLFszLDUsImJeey0xfV97XFxjVn1cXGNvbG9uICgxLGJeey0xfV97eSx4fSwxKSIsMl0sWzIsNSwiMV97XFxjVn0gXFxjb2xvbiAoMSwxLDEpIl0sWzYsNywiYiIsMCx7InNob3J0ZW4iOnsic291cmNlIjoyMCwidGFyZ2V0IjoyMH19XSxbMTAsOSwiYl57LTF9IiwwLHsic2hvcnRlbiI6eyJzb3VyY2UiOjIwLCJ0YXJnZXQiOjIwfX1dXQ==
\[\begin{tikzcd}[ampersand replacement=\&,cramped]
	\& {(\cV,b_{x,y})} \&\& {\BrMonCat_l} \\
	\\
	{(r(\cV),b_{y,x})} \&\&\& {\BrMonCat_l} \\
	\\
	\& {(\cV,b_{x,y})} \&\& {\BrMonCat_{l}}
	\arrow["{r(b_{\cV})\colon(1,b_{y,x},1)}"', from=1-2, to=3-1]
	\arrow["{1_{\cV} \colon (1,1,1)}", from=1-2, to=5-2]
	\arrow[""{name=0, anchor=center, inner sep=0}, "r", shift left=3, curve={height=-6pt}, from=1-4, to=3-4]
	\arrow[""{name=1, anchor=center, inner sep=0}, "1"', shift right=3, curve={height=6pt}, from=1-4, to=3-4]
	\arrow["{b^{-1}_{\cV}\colon (1,b^{-1}_{y,x},1)}"', from=3-1, to=5-2]
	\arrow[""{name=2, anchor=center, inner sep=0}, "1"', shift right=3, curve={height=6pt}, from=3-4, to=5-4]
	\arrow[""{name=3, anchor=center, inner sep=0}, "r", shift left=3, curve={height=-6pt}, from=3-4, to=5-4]
	\arrow["b", between={0.2}{0.8}, Rightarrow, nfold, from=0, to=1]
	\arrow["{b^{-1}}", between={0.2}{0.8}, Rightarrow, nfold, from=3, to=2]
\end{tikzcd}\cdpunct[4pt]{.}\]

Expectedly, if $\cV$ is a symmetric monoidal category then both the $2$-cells $b$ and $b^{-1}$ are identical and become part of a morphism of involutions. In general, we can compose this morphism of inverses $(b,b^{-1})$ to define other variants of the reversal duality $2$-cell, such as:
% https://q.uiver.app/#q=WzAsOSxbMCwwLCJcXEJyTW9uQ2F0X2wiXSxbMCwyLCJcXEJyTW9uQ2F0X2wiXSxbMiwwLCJcXE1vbkNhdF9sIl0sWzIsMiwiXFxNb25DYXRfbCJdLFs0LDAsIlxcQ2F0Il0sWzQsMiwiXFxDYXQiXSxbMCwzLCJcXE1vbihcXGNWKSJdLFsyLDMsIlxcTW9uKHIoXFxjVikpIl0sWzQsMywiXFxNb24oXFxjVikiXSxbMCwxLCJyIl0sWzAsMSwiMSIsMix7Im9mZnNldCI6MywiY3VydmUiOjN9XSxbNCw1LCIxIiwyXSxbMiw0LCJcXE1vbigtKSJdLFswLDIsIlUiXSxbMSwzLCJVIiwyXSxbMyw1LCJcXE1vbigtKSIsMl0sWzIsMywiciIsMl0sWzQsMywiciIsMix7InNob3J0ZW4iOnsic291cmNlIjoyMCwidGFyZ2V0IjoyMH0sImxldmVsIjoyfV0sWzYsNywicl97XFxjVn0iXSxbNyw4LCJcXE1vbihiX3tcXGNWfSkiXSxbOSwxMCwiYiIsMix7InNob3J0ZW4iOnsic291cmNlIjoyMCwidGFyZ2V0IjoyMH19XV0=
\[\begin{tikzcd}[ampersand replacement=\&]
	{\BrMonCat_l} \&\& {\MonCat_l} \&\& \Cat \\
	\\
	{\BrMonCat_l} \&\& {\MonCat_l} \&\& \Cat \\
	{\Mon(\cV)} \&\& {\Mon(r(\cV))} \&\& {\Mon(\cV)}
	\arrow["U", from=1-1, to=1-3]
	\arrow[""{name=0, anchor=center, inner sep=0}, "r", from=1-1, to=3-1]
	\arrow[""{name=1, anchor=center, inner sep=0}, "1"', shift right=3, curve={height=18pt}, from=1-1, to=3-1]
	\arrow["{\Mon(-)}", from=1-3, to=1-5]
	\arrow["r"', from=1-3, to=3-3]
	\arrow["r"', between={0.2}{0.8}, Rightarrow, nfold, from=1-5, to=3-3]
	\arrow["1"', from=1-5, to=3-5]
	\arrow["U"', from=3-1, to=3-3]
	\arrow["{\Mon(-)}"', from=3-3, to=3-5]
	\arrow["{r_{\cV}}", from=4-1, to=4-3]
	\arrow["{\Mon(b_{\cV})}", from=4-3, to=4-5]
	\arrow["b"', between={0.2}{0.8}, Rightarrow, nfold, from=0, to=1]
\end{tikzcd}\cdpunct[8pt]{.}\]
To illustrate the distinction between these variants of reversal, we provide the following example.

\begin{example}
	Let $\cV$ be the symmetric monoidal category $\Set$ with the symmetry given by the swap function $s_{x,y}: x \times y \rightarrow y \times x$. In this case, the component $b_{\Set}$ is given by the braided monoidal functor $(1_{\Set},s_{x,y},1) : (r(\Set),s_{y,x}) \rightarrow (\Set,s_{x,y})$. The morphism $r_{\Set}$ merely reinterprets a given monoid $M \in \Mon(\Set)$ as a monoid in $\Mon(r(\Set))$ in the sense depicted by
	% https://q.uiver.app/#q=WzAsNSxbMCwwLCJNIFxcdGltZXMgTSJdLFswLDIsIk0iXSxbMSwxLCI9Il0sWzIsMCwiTSBcXHRpbWVzX3tyfSBNIl0sWzIsMiwiTSJdLFswLDEsIlxcbXUiLDJdLFszLDQsIlxcbXUiXV0=
	\[\begin{tikzcd}[ampersand replacement=\&]
		{M \times M} \&\& {M \times_{r} M} \\
		\& {=} \\
		M \&\& M
		\arrow["\mu"', from=1-1, to=3-1]
		\arrow["\mu", from=1-3, to=3-3]
	\end{tikzcd}\cdpunct[8pt]{.}\]
	Thus, we should view the morphism $r_{\Set}$ as a comprehension duality. In contrast, the morphism $\Mon(b_{\Set})$ then transfers this monoid back to $\Mon(\Set)$ by precomposing with the symmetry as in
	% https://q.uiver.app/#q=WzAsNSxbMCwxLCJNIFxcdGltZXNfe3J9IE0iXSxbMCwzLCJNIl0sWzIsMiwiTSBcXHRpbWVzX3tyfSBNIl0sWzIsNCwiTSJdLFsyLDAsIk0gXFx0aW1lcyBNIl0sWzAsMSwiXFxtdSIsMl0sWzIsMywiXFxtdSJdLFs0LDIsInNfe00sTX0iXSxbNSwyLCIiLDAseyJzaG9ydGVuIjp7InNvdXJjZSI6MjAsInRhcmdldCI6MjB9LCJsZXZlbCI6MSwic3R5bGUiOnsidGFpbCI6eyJuYW1lIjoibWFwcyB0byJ9fX1dXQ==
	\[\begin{tikzcd}[ampersand replacement=\&]
		\&\& {M \times M} \\
		{M \times_{r} M} \\
		\&\& {M \times_{r} M} \\
		M \\
		\&\& M
		\arrow["{s_{M,M}}", from=1-3, to=3-3]
		\arrow[""{name=0, anchor=center, inner sep=0}, "\mu"', from=2-1, to=4-1]
		\arrow["\mu", from=3-3, to=5-3]
		\arrow[between={0.2}{0.8}, maps to, from=0, to=3-3]
	\end{tikzcd}\cdpunct[8pt]{.}\]
	The key takeaway is that one could replace $\Mon$ with $\mathbf{PsMon}$ and $\Set$ with $\mathbf{\Cat}$ to draw similar observation about the duality $r : \MonCat \rightarrow \MonCat$ itself. We discuss this later in \cref{sub:highermonduals}.
\end{example}

\begin{comment}
	Finally, $\BrMonCat_l$ also allows for another duality through the inversion of braiding, which gives the duality $2$-cell
	% https://q.uiver.app/#q=WzAsNCxbMCwwLCJcXEJyTW9uQ2F0X2wiXSxbNCwwLCJcXENhdCJdLFs0LDIsIlxcQ2F0Il0sWzAsMiwiXFxCck1vbkNhdF9sIl0sWzAsMSwiXFxjTW9uKC0pIl0sWzEsMiwiMSJdLFswLDMsImkiLDJdLFszLDIsIlxcY01vbigtKSIsMl0sWzEsMywiaSIsMix7InNob3J0ZW4iOnsic291cmNlIjoyMCwidGFyZ2V0IjoyMH0sImxldmVsIjoyfV1d
	\[\begin{tikzcd}[ampersand replacement=\&]
		{\BrMonCat_l} \&\&\&\& \Cat \\
		\\
		{\BrMonCat_l} \&\&\&\& \Cat
		\arrow["{\cMon(-)}", from=1-1, to=1-5]
		\arrow["i"', from=1-1, to=3-1]
		\arrow["i"', between={0.2}{0.8}, Rightarrow, from=1-5, to=3-1]
		\arrow["1", from=1-5, to=3-5]
		\arrow["{\cMon(-)}"', from=3-1, to=3-5]
	\end{tikzcd}\cdpunct[8pt]{,}\]
	satisfying the relations of being a morphism of involution. However, this duality restricts to the trivial $2$-cell when restricted to symmetric monoidal categories and extends to the trivial $2$-cell under the forgetful functor to monoidal categories. Moreover, $\SymMonCat$ doesn't admit any new non-trivial $2$-cell duality unlike $\BrMonCat$, as already discussed.
	
	\begin{remark}
		All duality $2$-cell commutes with each other in the obvious ways and also commutes with restrictions to commutative variants.
	\end{remark}
\end{comment}

A similar theory of dualities can be developed for arbitrary bimodules in a monoidal category. In particular, we could define a $2$-functor $\Mod(-)$ that takes the input as a monoidal category and outputs the bicategory of bimodules over arbitrary monoids (as long as we restrict our attention to monoidally cocomplete categories). However, this is rather unwieldy and introduces unnecessary assumptions. A more natural framework to encode these dualities is that of virtual double categories. See \cite[Defs.~2.1, 3.1, and 7.6]{shulman2010vequip} for a complete account and also see \cite[Ch.~5]{leinster2004} for our conventions on horizontal $1$-cells. In particular, we shall generally read horiztonal $1$-cells $M \colon A \rightarrow B$ as  $(B,A)$-bimodule.

Note that in the original definition, one can vertically reverse the direction of the square multicell to obtain the dual notion of a \enquote{co}-virtual double category. As the corresponding definitions are apparent in this dual case, we shall not write them explicitly. 

The only non-trivial result which we requires is that $\vcoDblCat$ and $\vDblCat$ exhibits a vertical duality involution in the following sense.

\begin{lemma}
	There exist two $2$-functors \[\op_{V}\colon \op_2(\vcoDblCat) \rightarrow \vDblCat\] and \[\op^t_{V} \colon \op_2(\vDblCat) \rightarrow \vcoDblCat,\]
	which form a morphism of inverses relative to the $\op_2$ involution. 
\end{lemma}

\begin{proof}
	we shall mainly define the action of $\op_V$ on each co-virtual double categories $\bV$. The rest is technical but straightforward. The data of $\op_{V}(\bV)$ is given by
	\begin{description}
		\item[Objects] 
		% https://q.uiver.app/#q=WzAsNCxbMCwxLCJYIl0sWzQsMSwiWCJdLFswLDAsIlxcYlYiXSxbNCwwLCJcXG9wX1YoXFxiVikiXV0=
		\[\begin{tikzcd}[ampersand replacement=\&,cramped]
			\bV \&\&\&\& {\op_V(\bV)} \\
			X \&\&\&\& X
		\end{tikzcd}\]
		\item[V $1$-cell] 
		% https://q.uiver.app/#q=WzAsNixbMCwzLCJZIl0sWzAsMSwiWCJdLFs0LDEsIlgiXSxbNCwzLCJZIl0sWzAsMCwiXFxiViJdLFs0LDAsIlxcb3BfVihcXGJWKSJdLFswLDFdLFsyLDNdXQ==
		\[\begin{tikzcd}[ampersand replacement=\&,cramped]
			\bV \&\&\&\& {\op_V(\bV)} \\
			X \&\&\&\& X \\
			\\
			Y \&\&\&\& Y
			\arrow[from=2-5, to=4-5]
			\arrow[from=4-1, to=2-1]
		\end{tikzcd}\]
		\item[H $1$-cell] 
		% https://q.uiver.app/#q=WzAsNixbMCwxLCJYIl0sWzIsMSwiWSJdLFs2LDEsIlgiXSxbOCwxLCJZIl0sWzEsMCwiXFxiViJdLFs3LDAsIlxcb3BfVihcXGJWKSJdLFswLDEsIiIsMCx7InN0eWxlIjp7ImJvZHkiOnsibmFtZSI6ImJhcnJlZCJ9fX1dLFsyLDMsIiIsMCx7InN0eWxlIjp7ImJvZHkiOnsibmFtZSI6ImJhcnJlZCJ9fX1dXQ==
		\[\begin{tikzcd}[ampersand replacement=\&,cramped,column sep=small]
			\& \bV \&\&\&\&\&\& {\op_V(\bV)} \& \\
			X \&\& Y \&\&\&\& X \&\& Y
			\arrow["\shortmid"{marking}, from=2-1, to=2-3]
			\arrow["\shortmid"{marking}, from=2-7, to=2-9]
		\end{tikzcd}\]
		\item[Sq multicell] 
		% https://q.uiver.app/#q=WzAsMTQsWzAsMywiWV8wIl0sWzQsMywiWV8xIl0sWzAsMSwiWF8wIl0sWzQsMSwiWF9uIl0sWzEsMSwiXFxjZG90cyJdLFszLDEsIlxcY2RvdHMiXSxbNiwxLCJYXzAiXSxbNiwzLCJZXzAiXSxbMTAsMywiWV8xIl0sWzcsMSwiXFxjZG90cyJdLFs5LDEsIlxcY2RvdHMiXSxbMTAsMSwiWF9uIl0sWzIsMCwiXFxiViJdLFs4LDAsIlxcb3BfVihcXGJWKSJdLFswLDEsIiIsMix7InN0eWxlIjp7ImJvZHkiOnsibmFtZSI6ImJhcnJlZCJ9fX1dLFswLDJdLFsxLDNdLFsyLDQsIiIsMCx7InN0eWxlIjp7ImJvZHkiOnsibmFtZSI6ImJhcnJlZCJ9fX1dLFs1LDMsIiIsMCx7InN0eWxlIjp7ImJvZHkiOnsibmFtZSI6ImJhcnJlZCJ9fX1dLFs0LDUsIiIsMSx7InN0eWxlIjp7ImJvZHkiOnsibmFtZSI6ImJhcnJlZCJ9fX1dLFs2LDddLFs3LDgsIiIsMix7InN0eWxlIjp7ImJvZHkiOnsibmFtZSI6ImJhcnJlZCJ9fX1dLFs5LDEwLCIiLDAseyJzdHlsZSI6eyJib2R5Ijp7Im5hbWUiOiJiYXJyZWQifX19XSxbMTEsOF0sWzYsOSwiIiwwLHsic3R5bGUiOnsiYm9keSI6eyJuYW1lIjoiYmFycmVkIn19fV0sWzEwLDExLCIiLDAseyJzdHlsZSI6eyJib2R5Ijp7Im5hbWUiOiJiYXJyZWQifX19XSxbMTQsMTksIiIsMSx7InNob3J0ZW4iOnsic291cmNlIjoyMCwidGFyZ2V0IjoyMH19XSxbMjIsMjEsIiIsMSx7InNob3J0ZW4iOnsic291cmNlIjoyMCwidGFyZ2V0IjoyMH19XV0=
		\[\begin{tikzcd}[ampersand replacement=\&,cramped,column sep=small]
			\&\& \bV \&\&\&\&\&\& {\op_V(\bV)} \&\& \\
			{X_0} \& \cdots \&\& \cdots \& {X_n} \&\& {X_0} \& \cdots \&\& \cdots \& {X_n} \\
			\\
			{Y_0} \&\&\&\& {Y_1} \&\& {Y_0} \&\&\&\& {Y_1}
			\arrow["\shortmid"{marking}, from=2-1, to=2-2]
			\arrow[""{name=0, anchor=center, inner sep=0}, "\shortmid"{marking}, from=2-2, to=2-4]
			\arrow["\shortmid"{marking}, from=2-4, to=2-5]
			\arrow["\shortmid"{marking}, from=2-7, to=2-8]
			\arrow[from=2-7, to=4-7]
			\arrow[""{name=1, anchor=center, inner sep=0}, "\shortmid"{marking}, from=2-8, to=2-10]
			\arrow["\shortmid"{marking}, from=2-10, to=2-11]
			\arrow[from=2-11, to=4-11]
			\arrow[from=4-1, to=2-1]
			\arrow[""{name=2, anchor=center, inner sep=0}, "\shortmid"{marking}, from=4-1, to=4-5]
			\arrow[from=4-5, to=2-5]
			\arrow[""{name=3, anchor=center, inner sep=0}, "\shortmid"{marking}, from=4-7, to=4-11]
			\arrow[between={0.2}{0.8}, Rightarrow, nfold, from=1, to=3]
			\arrow[between={0.2}{0.8}, Rightarrow, nfold, from=2, to=0]
		\end{tikzcd}\]
	\end{description}
	All of the data on the right is formal data similar to how the category $\cC^{\op}$ is defined when given a category $\cC$. Therefore, we simply induce the compositions and units of $\op_V(\bV)$ from $\bV$ itself in an analogous manner. 
	
	More technically, if the underlying fc-graph of a virtual double category $\bV$ is given by $T(C_0) \xleftarrow{s} C_1 \xrightarrow{t} C_0$, we can obtain a dual co-fc-graph given by $T(C_0) \xleftarrow{t'} C_1 \xrightarrow{s'} C_0$ with $t'=s$ and $s'=t$. 
	
\end{proof}

We briefly recall that there exists a chain of embeddings $\MonCat_{l} \hookrightarrow (\BiCat_{2})_{l} \hookrightarrow \vDblCat$. The first one is well known and simply identifies one object bicategories with monoidal categories \cite{lack2008icons}. The second one, we can describe as follows:
% https://q.uiver.app/#q=WzAsMTEsWzEsMCwiQiJdLFs2LDAsIlxcbWF0aGJie0J9Il0sWzAsMiwieF8wIl0sWzIsMiwieF9uIl0sWzQsMSwieF8wIl0sWzQsMywieF8wIl0sWzUsMSwieF8xIl0sWzYsMSwieF8yIl0sWzcsMSwiXFxjZG90cyJdLFs4LDEsInhfbiJdLFs4LDMsInhfbiJdLFsyLDMsImUoZl9uICxmX3tuLTF9LCBcXGRvdHMgLGZfMSkiLDAseyJjdXJ2ZSI6LTN9XSxbMiwzLCJnXzEiLDIseyJjdXJ2ZSI6M31dLFs0LDUsIlxcdmFyZXBzaWxvbl97eF8wfSIsMl0sWzUsMTAsImdfMSIsMix7InN0eWxlIjp7ImJvZHkiOnsibmFtZSI6ImJhcnJlZCJ9fX1dLFs0LDYsImZfMSIsMCx7InN0eWxlIjp7ImJvZHkiOnsibmFtZSI6ImJhcnJlZCJ9fX1dLFs2LDcsImZfMiIsMCx7InN0eWxlIjp7ImJvZHkiOnsibmFtZSI6ImJhcnJlZCJ9fX1dLFs3LDgsIiIsMCx7InN0eWxlIjp7ImJvZHkiOnsibmFtZSI6ImJhcnJlZCJ9fX1dLFs4LDksIiIsMCx7InN0eWxlIjp7ImJvZHkiOnsibmFtZSI6ImJhcnJlZCJ9fX1dLFs5LDEwLCJcXHZhcmVwc2lsb25fe3hfbn0iXSxbMCwxLCIiLDIseyJzdHlsZSI6eyJ0YWlsIjp7Im5hbWUiOiJtYXBzIHRvIn19fV0sWzExLDEyLCJcXGFscGhhIiwwLHsic2hvcnRlbiI6eyJzb3VyY2UiOjIwLCJ0YXJnZXQiOjIwfX1dLFs3LDE0LCJcXGFscGhhIFxcY29sb24gKGZfblxcY2RvdHMgKWZfMilmXzEgXFxSaWdodGFycm93IGdfMSIsMSx7InNob3J0ZW4iOnsic291cmNlIjoyMCwidGFyZ2V0IjoyMH19XV0=
\[\begin{tikzcd}[ampersand replacement=\&]
	\& B \&\&\&\&\& {\mathbb{B}} \&\& \\
	\&\&\&\& {x_0} \& {x_1} \& {x_2} \& \cdots \& {x_n} \\
	{x_0} \&\& {x_n} \\
	\&\&\&\& {x_0} \&\&\&\& {x_n}
	\arrow[maps to, from=1-2, to=1-7]
	\arrow["{f_1}"{inner sep=.8ex}, "\shortmid"{marking}, from=2-5, to=2-6]
	\arrow["{\varepsilon_{x_0}}"', from=2-5, to=4-5]
	\arrow["{f_2}"{inner sep=.8ex}, "\shortmid"{marking}, from=2-6, to=2-7]
	\arrow["\shortmid"{marking}, from=2-7, to=2-8]
	\arrow["\shortmid"{marking}, from=2-8, to=2-9]
	\arrow["{\varepsilon_{x_n}}", from=2-9, to=4-9]
	\arrow[""{name=0, anchor=center, inner sep=0}, "{e(f_n ,f_{n-1}, \dots ,f_1)}", curve={height=-18pt}, from=3-1, to=3-3]
	\arrow[""{name=1, anchor=center, inner sep=0}, "{g_1}"', curve={height=18pt}, from=3-1, to=3-3]
	\arrow[""{name=2, anchor=center, inner sep=0}, "{g_1}"'{inner sep=.8ex}, "\shortmid"{marking}, from=4-5, to=4-9]
	\arrow["{\alpha \colon (f_n\cdots )f_2)f_1 \Rightarrow g_1}"{description}, between={0.2}{0.8}, Rightarrow, nfold, from=2-7, to=2]
	\arrow["\alpha", between={0.2}{0.8}, Rightarrow, nfold, from=0, to=1]
\end{tikzcd}\cdpunct[8pt]{.}\]

An analogous embedding $\MonCat_{c} \hookrightarrow (\BiCat_{2})_{c} \hookrightarrow \vcoDblCat$ holds for the colax case. The first one is again clear. For the second one, we have the similar following description:
% https://q.uiver.app/#q=WzAsMTEsWzEsMCwiQiJdLFs2LDAsIlxcbWF0aGJie0J9Il0sWzAsMiwieF8wIl0sWzIsMiwieF9uIl0sWzQsMSwieF8wIl0sWzQsMywieF8wIl0sWzUsMSwieF8xIl0sWzYsMSwieF8yIl0sWzcsMSwiXFxjZG90cyJdLFs4LDEsInhfbiJdLFs4LDMsInhfbiJdLFsyLDMsImUoZl9uICxmX3tuLTF9LCBcXGRvdHMgLGZfMSkiLDAseyJjdXJ2ZSI6LTN9XSxbMiwzLCJnXzEiLDIseyJjdXJ2ZSI6M31dLFs1LDQsIlxcdmFyZXBzaWxvbl97eF8wfSJdLFs1LDEwLCJnXzEiLDIseyJzdHlsZSI6eyJib2R5Ijp7Im5hbWUiOiJiYXJyZWQifX19XSxbNCw2LCJmXzEiLDAseyJzdHlsZSI6eyJib2R5Ijp7Im5hbWUiOiJiYXJyZWQifX19XSxbNiw3LCJmXzIiLDAseyJzdHlsZSI6eyJib2R5Ijp7Im5hbWUiOiJiYXJyZWQifX19XSxbNyw4LCIiLDAseyJzdHlsZSI6eyJib2R5Ijp7Im5hbWUiOiJiYXJyZWQifX19XSxbOCw5LCIiLDAseyJzdHlsZSI6eyJib2R5Ijp7Im5hbWUiOiJiYXJyZWQifX19XSxbMTAsOSwiXFx2YXJlcHNpbG9uX3t4X259IiwyXSxbMCwxLCIiLDIseyJzdHlsZSI6eyJ0YWlsIjp7Im5hbWUiOiJtYXBzIHRvIn19fV0sWzEyLDExLCJcXGFscGhhIiwyLHsic2hvcnRlbiI6eyJzb3VyY2UiOjIwLCJ0YXJnZXQiOjIwfX1dLFsxNCw3LCJcXGFscGhhIFxcY29sb24gZ18xIFxcUmlnaHRhcnJvdyAoZl9uXFxjZG90cyApZl8yKWZfMSIsMSx7InNob3J0ZW4iOnsic291cmNlIjoyMCwidGFyZ2V0IjoyMH19XV0=
\[\begin{tikzcd}[ampersand replacement=\&]
	\& B \&\&\&\&\& {\mathbb{B}} \&\& \\
	\&\&\&\& {x_0} \& {x_1} \& {x_2} \& \cdots \& {x_n} \\
	{x_0} \&\& {x_n} \\
	\&\&\&\& {x_0} \&\&\&\& {x_n}
	\arrow[maps to, from=1-2, to=1-7]
	\arrow["{f_1}"{inner sep=.8ex}, "\shortmid"{marking}, from=2-5, to=2-6]
	\arrow["{f_2}"{inner sep=.8ex}, "\shortmid"{marking}, from=2-6, to=2-7]
	\arrow["\shortmid"{marking}, from=2-7, to=2-8]
	\arrow["\shortmid"{marking}, from=2-8, to=2-9]
	\arrow[""{name=0, anchor=center, inner sep=0}, "{e(f_n ,f_{n-1}, \dots ,f_1)}", curve={height=-18pt}, from=3-1, to=3-3]
	\arrow[""{name=1, anchor=center, inner sep=0}, "{g_1}"', curve={height=18pt}, from=3-1, to=3-3]
	\arrow["{\varepsilon_{x_0}}", from=4-5, to=2-5]
	\arrow[""{name=2, anchor=center, inner sep=0}, "{g_1}"'{inner sep=.8ex}, "\shortmid"{marking}, from=4-5, to=4-9]
	\arrow["{\varepsilon_{x_n}}"', from=4-9, to=2-9]
	\arrow["\alpha"', between={0.2}{0.8}, Rightarrow, nfold, from=1, to=0]
	\arrow["{\alpha \colon g_1 \Rightarrow (f_n\cdots )f_2)f_1}"{description}, between={0.2}{0.8}, Rightarrow, nfold, from=2, to=2-7]
\end{tikzcd}\cdpunct[8pt]{.}\]

The important fact about these two embeddings is that they canonically extend to morphisms of involutions and morphisms of inverses. 

\begin{lemma}\leavevmode
	\begin{enumerate} 
		\item There exist morphisms of inverses whose one direction is given by
		% https://q.uiver.app/#q=WzAsNixbMCwwLCJcXG9wXzIoXFxNb25DYXRfe2N9KSJdLFsyLDAsIlxcb3BfMihcXEJpQ2F0X3syfSlfe2N9Il0sWzQsMCwiXFxvcF8yKFxcdmNvRGJsQ2F0KSJdLFs0LDIsIlxcdkRibENhdCJdLFsyLDIsIihcXEJpQ2F0XzIpX2wiXSxbMCwyLCJcXE1vbkNhdF97bH0iXSxbMCwxLCIiLDAseyJzdHlsZSI6eyJ0YWlsIjp7Im5hbWUiOiJob29rIiwic2lkZSI6InRvcCJ9fX1dLFsxLDIsIiIsMCx7InN0eWxlIjp7InRhaWwiOnsibmFtZSI6Imhvb2siLCJzaWRlIjoidG9wIn19fV0sWzIsMywiXFxvcF97Vn0iXSxbMSw0LCJcXG9wX3syfSJdLFswLDUsIlxcb3BfezF9IiwyXSxbNCwzLCIiLDEseyJzdHlsZSI6eyJ0YWlsIjp7Im5hbWUiOiJob29rIiwic2lkZSI6InRvcCJ9fX1dLFs1LDQsIiIsMSx7InN0eWxlIjp7InRhaWwiOnsibmFtZSI6Imhvb2siLCJzaWRlIjoidG9wIn19fV0sWzIsNCwiXFxjb25nIiwxLHsic2hvcnRlbiI6eyJzb3VyY2UiOjIwLCJ0YXJnZXQiOjIwfSwibGV2ZWwiOjJ9XV0=
		\[\begin{tikzcd}[ampersand replacement=\&]
			{\op_2(\MonCat_{c})} \&\& {\op_2(\BiCat_{2})_{c}} \&\& {\op_2(\vcoDblCat)} \\
			\\
			{\MonCat_{l}} \&\& {(\BiCat_2)_l} \&\& \vDblCat
			\arrow[hook, from=1-1, to=1-3]
			\arrow["{\op_{1}}"', from=1-1, to=3-1]
			\arrow[hook, from=1-3, to=1-5]
			\arrow["{\op_{2}}", from=1-3, to=3-3]
			\arrow["\cong"{description}, between={0.2}{0.8}, Rightarrow, nfold, from=1-5, to=3-3]
			\arrow["{\op_{V}}", from=1-5, to=3-5]
			\arrow[hook, from=3-1, to=3-3]
			\arrow[hook, from=3-3, to=3-5]
		\end{tikzcd}\cdpunct[8pt]{,}\]
		and whose dual generalized inverse is given similarly.
		
		\item There exist morphisms of involution as given by
		% https://q.uiver.app/#q=WzAsNixbNCwwLCJcXHZEYmxDYXQiXSxbMiwwLCIoXFxCaUNhdF8yKV9sIl0sWzAsMCwiXFxNb25DYXRfe2x9Il0sWzQsMiwiXFx2RGJsQ2F0Il0sWzIsMiwiKFxcQmlDYXRfMilfbCJdLFswLDIsIlxcTW9uQ2F0X3tsfSJdLFsxLDAsIiIsMSx7InN0eWxlIjp7InRhaWwiOnsibmFtZSI6Imhvb2siLCJzaWRlIjoidG9wIn19fV0sWzIsMSwiIiwxLHsic3R5bGUiOnsidGFpbCI6eyJuYW1lIjoiaG9vayIsInNpZGUiOiJ0b3AifX19XSxbMCwzLCJcXG9wX3tIfSJdLFsxLDQsIlxcb3BfMSIsMl0sWzIsNSwiciIsMl0sWzUsNCwiIiwyLHsic3R5bGUiOnsidGFpbCI6eyJuYW1lIjoiaG9vayIsInNpZGUiOiJ0b3AifX19XSxbNCwzLCIiLDIseyJzdHlsZSI6eyJ0YWlsIjp7Im5hbWUiOiJob29rIiwic2lkZSI6InRvcCJ9fX1dLFswLDQsIlxcY29uZyIsMSx7InNob3J0ZW4iOnsic291cmNlIjoyMCwidGFyZ2V0IjoyMH0sImxldmVsIjoyfV1d
		\[\begin{tikzcd}[ampersand replacement=\&]
			{\MonCat_{l}} \&\& {(\BiCat_2)_l} \&\& \vDblCat \\
			\\
			{\MonCat_{l}} \&\& {(\BiCat_2)_l} \&\& \vDblCat
			\arrow[hook, from=1-1, to=1-3]
			\arrow["r"', from=1-1, to=3-1]
			\arrow[hook, from=1-3, to=1-5]
			\arrow["{\op_1}"', from=1-3, to=3-3]
			\arrow["\cong"{description}, between={0.2}{0.8}, Rightarrow, nfold, from=1-5, to=3-3]
			\arrow["{\op_{H}}", from=1-5, to=3-5]
			\arrow[hook, from=3-1, to=3-3]
			\arrow[hook, from=3-3, to=3-5]
		\end{tikzcd}\cdpunct[8pt]{.}\]
	\end{enumerate}
\end{lemma}

Note that the $2$-cell appearing in the involution case primarily arises from the fact that $((f_n \circ f_{n-1}) \circ \cdots ) \circ  f_1 \neq f_n \circ (f_{n-1} \circ (\cdots \circ f_1))$ in an arbitrary bicategory. To illustrate this issue further, consider a left-biased lax monoidal category $\cV$, that is, the associator need not an invertible $2$-cell. Applying $r$ to this lax monoidal category doesn't yield a left-biased lax monoidal category, as there is no $a^{-1}$ available. At best, we obtain a canonical right-biased lax monoidal category. All of this suggests that a more thorough approach would require generalized inverses $r \colon \MonCat_l^{Lb} \rightarrow \MonCat_l^{Rb}$ and $r^t\colon \MonCat_l^{Rb}  \rightarrow \MonCat_{l}^{Lb}$. However, we are content with our approach here. 

\begin{comment}
	Finally, we state the similar lemma for the case of inversion. 
	
	\begin{lemma}
		The morphism of involution $i$ extends to the virtual double categories case through the following strict $2$-transformation cells:
		% https://q.uiver.app/#q=WzAsNCxbMCwwLCJcXEJyTW9uQ2F0X2wiXSxbMiwwLCJcXHZEYmxDYXQiXSxbMCwyLCJcXEJyTW9uQ2F0X2wiXSxbMiwyLCJcXHZEYmxDYXQiXSxbMSwzLCIxIl0sWzAsMiwiaSIsMl0sWzIsMywiXFxjTW9kKC0pIiwyXSxbMCwxLCJcXGNNb2QoLSkiXSxbMSwyLCJpIiwwLHsic2hvcnRlbiI6eyJzb3VyY2UiOjIwLCJ0YXJnZXQiOjIwfSwibGV2ZWwiOjJ9XV0=
		\[\begin{tikzcd}[ampersand replacement=\&]
			{\BrMonCat_l} \&\& \vDblCat \\
			\\
			{\BrMonCat_l} \&\& \vDblCat
			\arrow["{\cMod(-)}", from=1-1, to=1-3]
			\arrow["i"', from=1-1, to=3-1]
			\arrow["i", between={0.2}{0.8}, Rightarrow, nfold, from=1-3, to=3-1]
			\arrow["1", from=1-3, to=3-3]
			\arrow["{\cMod(-)}"', from=3-1, to=3-3]
		\end{tikzcd}\cdpunct[8pt]{.}\]
		Furhermore, this morphism can be factored through the virtual equipments.
	\end{lemma}
\end{comment}

We have now established that monoidal categories are a special case of a virtual double categories, as well as co-virtual double categories. This suggests that the dualities of monoidal categories should also extend to the virtual double categorical case. We begin by characterization the $\op_1$ duality.

\begin{theorem}
	The $2$-functors $\Mod \colon \vDblCat \rightarrow \vDblCat$ and $\CoMod \colon \vcoDblCat \rightarrow \vcoDblCat$ extend to a $1$-morphism of relative inverses by the strict $2$-natural transformations
	% https://q.uiver.app/#q=WzAsOCxbMCwwLCJcXG9wXzIoXFx2Y29EYmxDYXQpIl0sWzAsMiwiXFx2RGJsQ2F0Il0sWzMsMiwiXFx2RGJsQ2F0Il0sWzMsMCwiXFxvcF8yKFxcdmNvRGJsQ2F0KSJdLFswLDMsIlxcb3BfMihcXHZEYmxDYXQpIl0sWzMsMywiXFxvcF8yKFxcdkRibENhdCkiXSxbMCw1LCJcXHZjb0RibENhdCJdLFszLDUsIlxcdmNvRGJsQ2F0Il0sWzAsMSwiXFxvcF9WIiwyXSxbMywyLCJcXG9wX1YiXSxbMSwyLCJcXE1vZCgtKSIsMl0sWzAsMywiXFxvcF8yKFxcQ29Nb2QoLSkpIl0sWzMsMSwiXFxvcCIsMix7InNob3J0ZW4iOnsic291cmNlIjoyMCwidGFyZ2V0IjoyMH0sImxldmVsIjoyfV0sWzYsNywiXFxDb01vZCgtKSIsMl0sWzQsNiwiXFxvcF97Vn1ee3R9IiwyXSxbNSw3LCJcXG9wX3tWfV57dH0iXSxbNCw1LCJcXG9wXzIoXFxNb2QoLSkpIl0sWzUsNiwiXFxvcF57dH0iLDIseyJzaG9ydGVuIjp7InNvdXJjZSI6MjAsInRhcmdldCI6MjB9LCJsZXZlbCI6Mn1dXQ==
	\[\begin{tikzcd}[ampersand replacement=\&]
		{\op_2(\vcoDblCat)} \&\&\& {\op_2(\vcoDblCat)} \\
		\\
		\vDblCat \&\&\& \vDblCat \\
		{\op_2(\vDblCat)} \&\&\& {\op_2(\vDblCat)} \\
		\\
		\vcoDblCat \&\&\& \vcoDblCat
		\arrow["{\op_2(\CoMod(-))}", from=1-1, to=1-4]
		\arrow["{\op_V}"', from=1-1, to=3-1]
		\arrow["\op"', between={0.2}{0.8}, Rightarrow, nfold, from=1-4, to=3-1]
		\arrow["{\op_V}", from=1-4, to=3-4]
		\arrow["{\Mod(-)}"', from=3-1, to=3-4]
		\arrow["{\op_2(\Mod(-))}", from=4-1, to=4-4]
		\arrow["{\op_{V}^{t}}"', from=4-1, to=6-1]
		\arrow["{\op^{t}}"', between={0.2}{0.8}, Rightarrow, nfold, from=4-4, to=6-1]
		\arrow["{\op_{V}^{t}}", from=4-4, to=6-4]
		\arrow["{\CoMod(-)}"', from=6-1, to=6-4]
	\end{tikzcd}\cdpunct[8pt]{.}\]
\end{theorem}

\begin{proof}
	As usual, we need to define the components $\op_{\bV} : \op_{V}(\CoMod(\bV)) \rightarrow \Mod(\op_{V}(\bV))$.
	\begin{description}
		\item[Objects] 
		% https://q.uiver.app/#q=WzAsMTgsWzEsMCwiQyBcXGNvbG9uZXFxIChcXGJ1bGxldCxDLFxcbXUsZSkiXSxbNSwwLCJDXntcXG9wfSBcXGNvbG9uZXFxIChcXGJ1bGxldCxDLCBcXG11XntcXG9wfSwgZV57XFxvcH0pIl0sWzAsMiwiXFxidWxsZXQiXSxbMCwxLCJcXGJ1bGxldCJdLFsxLDEsIlxcYnVsbGV0Il0sWzIsMSwiXFxidWxsZXQiXSxbMiwyLCJcXGJ1bGxldCJdLFs0LDEsIlxcYnVsbGV0Il0sWzQsMiwiXFxidWxsZXQiXSxbNSwxLCJcXGJ1bGxldCJdLFs2LDEsIlxcYnVsbGV0Il0sWzYsMiwiXFxidWxsZXQiXSxbMCw0LCJcXGJ1bGxldCJdLFsyLDQsIlxcYnVsbGV0Il0sWzEsMywiXFxidWxsZXQiXSxbNCw0LCJcXGJ1bGxldCJdLFs2LDQsIlxcYnVsbGV0Il0sWzUsMywiXFxidWxsZXQiXSxbMyw0LCJDIiwwLHsic3R5bGUiOnsiYm9keSI6eyJuYW1lIjoiYmFycmVkIn19fV0sWzQsNSwiQyIsMCx7InN0eWxlIjp7ImJvZHkiOnsibmFtZSI6ImJhcnJlZCJ9fX1dLFsyLDMsIiIsMCx7InN0eWxlIjp7ImhlYWQiOnsibmFtZSI6Im5vbmUifX19XSxbMiw2LCJDIiwyLHsic3R5bGUiOnsiYm9keSI6eyJuYW1lIjoiYmFycmVkIn19fV0sWzYsNSwiIiwyLHsic3R5bGUiOnsiaGVhZCI6eyJuYW1lIjoibm9uZSJ9fX1dLFs3LDksIkMiLDAseyJzdHlsZSI6eyJib2R5Ijp7Im5hbWUiOiJiYXJyZWQifX19XSxbOSwxMCwiQyIsMCx7InN0eWxlIjp7ImJvZHkiOnsibmFtZSI6ImJhcnJlZCJ9fX1dLFsxMCwxMSwiIiwwLHsic3R5bGUiOnsiaGVhZCI6eyJuYW1lIjoibm9uZSJ9fX1dLFs4LDExLCJDIiwyLHsic3R5bGUiOnsiYm9keSI6eyJuYW1lIjoiYmFycmVkIn19fV0sWzcsOCwiIiwyLHsic3R5bGUiOnsiaGVhZCI6eyJuYW1lIjoibm9uZSJ9fX1dLFsxMiwxMywiQyIsMix7InN0eWxlIjp7ImJvZHkiOnsibmFtZSI6ImJhcnJlZCJ9fX1dLFsxMiwxNCwiIiwyLHsic3R5bGUiOnsiaGVhZCI6eyJuYW1lIjoibm9uZSJ9fX1dLFsxMywxNCwiIiwwLHsic3R5bGUiOnsiaGVhZCI6eyJuYW1lIjoibm9uZSJ9fX1dLFsxNywxNSwiIiwyLHsic3R5bGUiOnsiaGVhZCI6eyJuYW1lIjoibm9uZSJ9fX1dLFsxNywxNiwiIiwwLHsic3R5bGUiOnsiaGVhZCI6eyJuYW1lIjoibm9uZSJ9fX1dLFsxNSwxNiwiQyIsMix7InN0eWxlIjp7ImJvZHkiOnsibmFtZSI6ImJhcnJlZCJ9fX1dLFswLDEsIiIsMCx7InNob3J0ZW4iOnsic291cmNlIjoyMCwidGFyZ2V0IjoyMH0sInN0eWxlIjp7InRhaWwiOnsibmFtZSI6Im1hcHMgdG8ifX19XSxbMjEsNCwiXFxtdSIsMCx7InNob3J0ZW4iOnsic291cmNlIjoyMH19XSxbOSwyNiwiXFxtdV57XFxvcH0iLDAseyJzaG9ydGVuIjp7InRhcmdldCI6MjB9fV0sWzI4LDE0LCJlIiwwLHsic2hvcnRlbiI6eyJzb3VyY2UiOjIwfX1dLFsxNywzMywiZV57XFxvcH0iLDAseyJzaG9ydGVuIjp7InNvdXJjZSI6MjAsInRhcmdldCI6MjB9fV1d
		\[\begin{tikzcd}[ampersand replacement=\&]
			\& {C \coloneqq (\bullet,C,\mu,e)} \&\&\&\& {C^{\op} \coloneqq (\bullet,C, \mu^{\op}, e^{\op})} \& \\
			\bullet \& \bullet \& \bullet \&\& \bullet \& \bullet \& \bullet \\
			\bullet \&\& \bullet \&\& \bullet \&\& \bullet \\
			\& \bullet \&\&\&\& \bullet \\
			\bullet \&\& \bullet \&\& \bullet \&\& \bullet
			\arrow[between={0.2}{0.8}, maps to, from=1-2, to=1-6]
			\arrow["C"{inner sep=.8ex}, "\shortmid"{marking}, from=2-1, to=2-2]
			\arrow["C"{inner sep=.8ex}, "\shortmid"{marking}, from=2-2, to=2-3]
			\arrow["C"{inner sep=.8ex}, "\shortmid"{marking}, from=2-5, to=2-6]
			\arrow[no head, from=2-5, to=3-5]
			\arrow["C"{inner sep=.8ex}, "\shortmid"{marking}, from=2-6, to=2-7]
			\arrow[no head, from=2-7, to=3-7]
			\arrow[no head, from=3-1, to=2-1]
			\arrow[""{name=0, anchor=center, inner sep=0}, "C"'{inner sep=.8ex}, "\shortmid"{marking}, from=3-1, to=3-3]
			\arrow[no head, from=3-3, to=2-3]
			\arrow[""{name=1, anchor=center, inner sep=0}, "C"'{inner sep=.8ex}, "\shortmid"{marking}, from=3-5, to=3-7]
			\arrow[no head, from=4-6, to=5-5]
			\arrow[no head, from=4-6, to=5-7]
			\arrow[no head, from=5-1, to=4-2]
			\arrow[""{name=2, anchor=center, inner sep=0}, "C"'{inner sep=.8ex}, "\shortmid"{marking}, from=5-1, to=5-3]
			\arrow[no head, from=5-3, to=4-2]
			\arrow[""{name=3, anchor=center, inner sep=0}, "C"'{inner sep=.8ex}, "\shortmid"{marking}, from=5-5, to=5-7]
			\arrow["{\mu^{\op}}", between={0}{0.8}, Rightarrow, nfold, from=2-6, to=1]
			\arrow["\mu", between={0.2}{1}, Rightarrow, nfold, from=0, to=2-2]
			\arrow["{e^{\op}}", between={0.2}{0.8}, Rightarrow, nfold, from=4-6, to=3]
			\arrow["e", between={0.2}{1}, Rightarrow, nfold, from=2, to=4-2]
		\end{tikzcd}\]
		\item[V $1$-cell] 
		% https://q.uiver.app/#q=WzAsOCxbMCwwLCJcXGJ1bGxldCJdLFsyLDAsIlxcYnVsbGV0Il0sWzAsMiwiXFxzdGFyIl0sWzIsMiwiXFxzdGFyIl0sWzUsMCwiXFxidWxsZXQiXSxbNywwLCJcXGJ1bGxldCJdLFs1LDIsIlxcc3RhciJdLFs3LDIsIlxcc3RhciJdLFswLDEsIkMiLDAseyJzdHlsZSI6eyJib2R5Ijp7Im5hbWUiOiJiYXJyZWQifX19XSxbMiwwLCJmX3swfSJdLFszLDEsImZfezB9IiwyXSxbMiwzLCJEIiwyLHsic3R5bGUiOnsiYm9keSI6eyJuYW1lIjoiYmFycmVkIn19fV0sWzYsNywiRF57XFxvcH0iLDIseyJzdHlsZSI6eyJib2R5Ijp7Im5hbWUiOiJiYXJyZWQifX19XSxbNCw1LCJDXntcXG9wfSIsMCx7InN0eWxlIjp7ImJvZHkiOnsibmFtZSI6ImJhcnJlZCJ9fX1dLFs0LDYsImZfMF57XFxvcH0iLDJdLFs1LDcsImZfMF57XFxvcH0iXSxbMTEsOCwiZiIsMix7InNob3J0ZW4iOnsic291cmNlIjoyMCwidGFyZ2V0IjoyMH19XSxbMTMsMTIsImZee1xcb3B9IiwwLHsic2hvcnRlbiI6eyJzb3VyY2UiOjIwLCJ0YXJnZXQiOjIwfX1dLFsxMCwxNCwiIiwyLHsic2hvcnRlbiI6eyJzb3VyY2UiOjIwLCJ0YXJnZXQiOjIwfSwibGV2ZWwiOjEsInN0eWxlIjp7InRhaWwiOnsibmFtZSI6Im1hcHMgdG8ifX19XV0=
		\[\begin{tikzcd}[ampersand replacement=\&]
			\bullet \&\& \bullet \&\&\& \bullet \&\& \bullet \\
			\\
			\star \&\& \star \&\&\& \star \&\& \star
			\arrow[""{name=0, anchor=center, inner sep=0}, "C"{inner sep=.8ex}, "\shortmid"{marking}, from=1-1, to=1-3]
			\arrow[""{name=1, anchor=center, inner sep=0}, "{C^{\op}}"{inner sep=.8ex}, "\shortmid"{marking}, from=1-6, to=1-8]
			\arrow[""{name=2, anchor=center, inner sep=0}, "{f_0^{\op}}"', from=1-6, to=3-6]
			\arrow["{f_0^{\op}}", from=1-8, to=3-8]
			\arrow["{f_{0}}", from=3-1, to=1-1]
			\arrow[""{name=3, anchor=center, inner sep=0}, "D"'{inner sep=.8ex}, "\shortmid"{marking}, from=3-1, to=3-3]
			\arrow[""{name=4, anchor=center, inner sep=0}, "{f_{0}}"', from=3-3, to=1-3]
			\arrow[""{name=5, anchor=center, inner sep=0}, "{D^{\op}}"'{inner sep=.8ex}, "\shortmid"{marking}, from=3-6, to=3-8]
			\arrow["{f^{\op}}", between={0.2}{0.8}, Rightarrow, nfold, from=1, to=5]
			\arrow["f"', between={0.2}{0.8}, Rightarrow, nfold, from=3, to=0]
			\arrow[between={0.2}{0.8}, maps to, from=4, to=2]
		\end{tikzcd}\]
		\item[H $1$-cell]
		% https://q.uiver.app/#q=WzAsMjgsWzQsMSwiXFxzdGFyIl0sWzAsMywiXFxidWxsZXQiXSxbMCwyLCJcXGJ1bGxldCJdLFsyLDIsIlxcYnVsbGV0Il0sWzQsMiwiXFxzdGFyIl0sWzQsMywiXFxzdGFyIl0sWzAsNCwiXFxidWxsZXQiXSxbMCw1LCJcXGJ1bGxldCJdLFsyLDQsIlxcc3RhciJdLFs0LDQsIlxcc3RhciJdLFs0LDUsIlxcc3RhciJdLFswLDEsIlxcYnVsbGV0Il0sWzYsMSwiXFxidWxsZXQiXSxbMTAsMSwiXFxzdGFyIl0sWzYsMiwiXFxidWxsZXQiXSxbOCwyLCJcXGJ1bGxldCJdLFs2LDMsIlxcYnVsbGV0Il0sWzgsNCwiXFxzdGFyIl0sWzEwLDQsIlxcc3RhciJdLFsxMCw1LCJcXHN0YXIiXSxbNiw1LCJcXGJ1bGxldCJdLFs2LDQsIlxcYnVsbGV0Il0sWzEwLDMsIlxcc3RhciJdLFsxMCwyLCJcXHN0YXIiXSxbMCwwLCIoQyxcXGJ1bGxldCkiXSxbNCwwLCIoRCxcXHN0YXIpIl0sWzYsMCwiKENee1xcb3B9LFxcYnVsbGV0KSJdLFsxMCwwLCIoRF57XFxvcH0sXFxzdGFyKSJdLFsyLDMsIkMiLDAseyJzdHlsZSI6eyJib2R5Ijp7Im5hbWUiOiJiYXJyZWQifX19XSxbMyw0LCJYIiwwLHsic3R5bGUiOnsiYm9keSI6eyJuYW1lIjoiYmFycmVkIn19fV0sWzEsMiwiIiwwLHsic3R5bGUiOnsiaGVhZCI6eyJuYW1lIjoibm9uZSJ9fX1dLFsxLDUsIlgiLDIseyJzdHlsZSI6eyJib2R5Ijp7Im5hbWUiOiJiYXJyZWQifX19XSxbNSw0LCIiLDIseyJzdHlsZSI6eyJoZWFkIjp7Im5hbWUiOiJub25lIn19fV0sWzYsOCwiWCIsMCx7InN0eWxlIjp7ImJvZHkiOnsibmFtZSI6ImJhcnJlZCJ9fX1dLFs4LDksIkQiLDAseyJzdHlsZSI6eyJib2R5Ijp7Im5hbWUiOiJiYXJyZWQifX19XSxbOSwxMCwiIiwwLHsic3R5bGUiOnsiaGVhZCI6eyJuYW1lIjoibm9uZSJ9fX1dLFs3LDEwLCJYIiwyLHsic3R5bGUiOnsiYm9keSI6eyJuYW1lIjoiYmFycmVkIn19fV0sWzYsNywiIiwyLHsic3R5bGUiOnsiaGVhZCI6eyJuYW1lIjoibm9uZSJ9fX1dLFsxMSwwLCJYIiwwLHsic3R5bGUiOnsiYm9keSI6eyJuYW1lIjoiYmFycmVkIn19fV0sWzEyLDEzLCJYIiwwLHsic3R5bGUiOnsiYm9keSI6eyJuYW1lIjoiYmFycmVkIn19fV0sWzE2LDIyLCJYXntcXG9wfSIsMix7InN0eWxlIjp7ImJvZHkiOnsibmFtZSI6ImJhcnJlZCJ9fX1dLFsxNCwxNSwiQ157XFxvcH0iLDAseyJzdHlsZSI6eyJib2R5Ijp7Im5hbWUiOiJiYXJyZWQifX19XSxbMjEsMTcsIlhee1xcb3B9IiwwLHsic3R5bGUiOnsiYm9keSI6eyJuYW1lIjoiYmFycmVkIn19fV0sWzE3LDE4LCJEXntcXG9wfSIsMCx7InN0eWxlIjp7ImJvZHkiOnsibmFtZSI6ImJhcnJlZCJ9fX1dLFsyMCwxOSwiWF57XFxvcH0iLDIseyJzdHlsZSI6eyJib2R5Ijp7Im5hbWUiOiJiYXJyZWQifX19XSxbMTUsMjMsIlhee1xcb3B9IiwwLHsic3R5bGUiOnsiYm9keSI6eyJuYW1lIjoiYmFycmVkIn19fV0sWzE0LDE2LCIiLDEseyJzdHlsZSI6eyJoZWFkIjp7Im5hbWUiOiJub25lIn19fV0sWzIzLDIyLCIiLDEseyJzdHlsZSI6eyJoZWFkIjp7Im5hbWUiOiJub25lIn19fV0sWzE4LDE5LCIiLDEseyJzdHlsZSI6eyJoZWFkIjp7Im5hbWUiOiJub25lIn19fV0sWzIxLDIwLCIiLDEseyJzdHlsZSI6eyJoZWFkIjp7Im5hbWUiOiJub25lIn19fV0sWzI2LDI3LCJYXntcXG9wfVxcY29sb25lcXEoWCxcXGxhbWJkYV57XFxvcH0sXFxyaG9ee1xcb3B9KSIsMCx7InN0eWxlIjp7ImJvZHkiOnsibmFtZSI6ImJhcnJlZCJ9fX1dLFsyNSwyNiwiIiwwLHsic2hvcnRlbiI6eyJzb3VyY2UiOjIwLCJ0YXJnZXQiOjIwfSwic3R5bGUiOnsidGFpbCI6eyJuYW1lIjoibWFwcyB0byJ9fX1dLFsyNCwyNSwiWCBcXGNvbG9uZXFxKFgsXFxsYW1iZGEsXFxyaG8pIiwwLHsic3R5bGUiOnsiYm9keSI6eyJuYW1lIjoiYmFycmVkIn19fV0sWzMxLDMsIlxccmhvIiwwLHsic2hvcnRlbiI6eyJzb3VyY2UiOjIwfX1dLFszNiw4LCJcXGxhbWJkYSIsMCx7InNob3J0ZW4iOnsic291cmNlIjoyMH19XSxbMTUsNDAsIlxccmhvXntcXG9wfSIsMix7InNob3J0ZW4iOnsidGFyZ2V0IjoyMH19XSxbMTcsNDQsIlxcbGFtYmRhXntcXG9wfSIsMix7InNob3J0ZW4iOnsidGFyZ2V0IjoyMH19XV0=
		\[\begin{tikzcd}[ampersand replacement=\&,column sep=small]
			{(C,\bullet)} \&\&\&\& {(D,\star)} \&\& {(C^{\op},\bullet)} \&\&\&\& {(D^{\op},\star)} \\
			\bullet \&\&\&\& \star \&\& \bullet \&\&\&\& \star \\
			\bullet \&\& \bullet \&\& \star \&\& \bullet \&\& \bullet \&\& \star \\
			\bullet \&\&\&\& \star \&\& \bullet \&\&\&\& \star \\
			\bullet \&\& \star \&\& \star \&\& \bullet \&\& \star \&\& \star \\
			\bullet \&\&\&\& \star \&\& \bullet \&\&\&\& \star
			\arrow["{X \coloneqq(X,\lambda,\rho)}"{inner sep=.8ex}, "\shortmid"{marking}, from=1-1, to=1-5]
			\arrow[between={0.2}{0.8}, maps to, from=1-5, to=1-7]
			\arrow["{X^{\op}\coloneqq(X,\lambda^{\op},\rho^{\op})}"{inner sep=.8ex}, "\shortmid"{marking}, from=1-7, to=1-11]
			\arrow["X"{inner sep=.8ex}, "\shortmid"{marking}, from=2-1, to=2-5]
			\arrow["X"{inner sep=.8ex}, "\shortmid"{marking}, from=2-7, to=2-11]
			\arrow["C"{inner sep=.8ex}, "\shortmid"{marking}, from=3-1, to=3-3]
			\arrow["X"{inner sep=.8ex}, "\shortmid"{marking}, from=3-3, to=3-5]
			\arrow["{C^{\op}}"{inner sep=.8ex}, "\shortmid"{marking}, from=3-7, to=3-9]
			\arrow[no head, from=3-7, to=4-7]
			\arrow["{X^{\op}}"{inner sep=.8ex}, "\shortmid"{marking}, from=3-9, to=3-11]
			\arrow[no head, from=3-11, to=4-11]
			\arrow[no head, from=4-1, to=3-1]
			\arrow[""{name=0, anchor=center, inner sep=0}, "X"'{inner sep=.8ex}, "\shortmid"{marking}, from=4-1, to=4-5]
			\arrow[no head, from=4-5, to=3-5]
			\arrow[""{name=1, anchor=center, inner sep=0}, "{X^{\op}}"'{inner sep=.8ex}, "\shortmid"{marking}, from=4-7, to=4-11]
			\arrow["X"{inner sep=.8ex}, "\shortmid"{marking}, from=5-1, to=5-3]
			\arrow[no head, from=5-1, to=6-1]
			\arrow["D"{inner sep=.8ex}, "\shortmid"{marking}, from=5-3, to=5-5]
			\arrow[no head, from=5-5, to=6-5]
			\arrow["{X^{\op}}"{inner sep=.8ex}, "\shortmid"{marking}, from=5-7, to=5-9]
			\arrow[no head, from=5-7, to=6-7]
			\arrow["{D^{\op}}"{inner sep=.8ex}, "\shortmid"{marking}, from=5-9, to=5-11]
			\arrow[no head, from=5-11, to=6-11]
			\arrow[""{name=2, anchor=center, inner sep=0}, "X"'{inner sep=.8ex}, "\shortmid"{marking}, from=6-1, to=6-5]
			\arrow[""{name=3, anchor=center, inner sep=0}, "{X^{\op}}"'{inner sep=.8ex}, "\shortmid"{marking}, from=6-7, to=6-11]
			\arrow["{\rho^{\op}}"', between={0}{0.8}, Rightarrow, nfold, from=3-9, to=1]
			\arrow["\rho", between={0.2}{1}, Rightarrow, nfold, from=0, to=3-3]
			\arrow["{\lambda^{\op}}"', between={0}{0.8}, Rightarrow, nfold, from=5-9, to=3]
			\arrow["\lambda", between={0.2}{1}, Rightarrow, nfold, from=2, to=5-3]
		\end{tikzcd}\]
		\item[Sq multicell]
		% https://q.uiver.app/#q=WzAsMTAsWzAsMiwiXFxidWxsZXQiXSxbNCwyLCJcXHN0YXIiXSxbMCwwLCJcXGJ1bGxldF8xIl0sWzQsMCwiXFxidWxsZXRfbiJdLFsyLDAsIlxcY2RvdHMiXSxbNiwwLCJcXGJ1bGxldF8xIl0sWzgsMCwiXFxjZG90cyJdLFsxMCwwLCJcXGJ1bGxldF9uIl0sWzYsMiwiXFxidWxsZXQiXSxbMTAsMiwiXFxzdGFyIl0sWzIsNCwiWF8xIiwwLHsic3R5bGUiOnsiYm9keSI6eyJuYW1lIjoiYmFycmVkIn19fV0sWzQsMywiWF9uIiwwLHsic3R5bGUiOnsiYm9keSI6eyJuYW1lIjoiYmFycmVkIn19fV0sWzAsMSwiWSIsMix7InN0eWxlIjp7ImJvZHkiOnsibmFtZSI6ImJhcnJlZCJ9fX1dLFswLDIsImYiXSxbMSwzLCJnIiwyXSxbNSw4LCJmXntcXG9wfSIsMl0sWzcsOSwiZ157XFxvcH0iXSxbOCw5LCJZXntcXG9wfSIsMix7InN0eWxlIjp7ImJvZHkiOnsibmFtZSI6ImJhcnJlZCJ9fX1dLFs1LDYsIlhfMV57XFxvcH0iLDAseyJzdHlsZSI6eyJib2R5Ijp7Im5hbWUiOiJiYXJyZWQifX19XSxbNiw3LCJYX25ee1xcb3B9IiwwLHsic3R5bGUiOnsiYm9keSI6eyJuYW1lIjoiYmFycmVkIn19fV0sWzEyLDQsIlxcYWxwaGEiLDIseyJzaG9ydGVuIjp7InNvdXJjZSI6MjB9fV0sWzYsMTcsIlxcYWxwaGFee1xcb3B9IiwwLHsic2hvcnRlbiI6eyJ0YXJnZXQiOjIwfX1dLFsxNCwxNSwiIiwyLHsic2hvcnRlbiI6eyJzb3VyY2UiOjMwLCJ0YXJnZXQiOjMwfSwibGV2ZWwiOjEsInN0eWxlIjp7InRhaWwiOnsibmFtZSI6Im1hcHMgdG8ifX19XV0=
		\[\begin{tikzcd}[ampersand replacement=\&]
			{\bullet_1} \&\& \cdots \&\& {\bullet_n} \&\& {\bullet_1} \&\& \cdots \&\& {\bullet_n} \\
			\\
			\bullet \&\&\&\& \star \&\& \bullet \&\&\&\& \star
			\arrow["{X_1}"{inner sep=.8ex}, "\shortmid"{marking}, from=1-1, to=1-3]
			\arrow["{X_n}"{inner sep=.8ex}, "\shortmid"{marking}, from=1-3, to=1-5]
			\arrow["{X_1^{\op}}"{inner sep=.8ex}, "\shortmid"{marking}, from=1-7, to=1-9]
			\arrow[""{name=0, anchor=center, inner sep=0}, "{f^{\op}}"', from=1-7, to=3-7]
			\arrow["{X_n^{\op}}"{inner sep=.8ex}, "\shortmid"{marking}, from=1-9, to=1-11]
			\arrow["{g^{\op}}", from=1-11, to=3-11]
			\arrow["f", from=3-1, to=1-1]
			\arrow[""{name=1, anchor=center, inner sep=0}, "Y"'{inner sep=.8ex}, "\shortmid"{marking}, from=3-1, to=3-5]
			\arrow[""{name=2, anchor=center, inner sep=0}, "g"', from=3-5, to=1-5]
			\arrow[""{name=3, anchor=center, inner sep=0}, "{Y^{\op}}"'{inner sep=.8ex}, "\shortmid"{marking}, from=3-7, to=3-11]
			\arrow["{\alpha^{\op}}", between={0}{0.8}, Rightarrow, nfold, from=1-9, to=3]
			\arrow["\alpha"', between={0.2}{1}, Rightarrow, nfold, from=1, to=1-3]
			\arrow[between={0.3}{0.7}, maps to, from=2, to=0]
		\end{tikzcd}\]
	\end{description}
	The other family of components can be defined similarly.
\end{proof}

The case for reversal duality is exactly the same as above.

\begin{theorem}
	The $2$-functor $\Mod \colon \vDblCat \rightarrow \vDblCat$ extends to a $1$-morphism of involutions by the strict $2$-natural transformation
	% https://q.uiver.app/#q=WzAsNCxbMCwwLCJcXHZEYmxDYXQiXSxbMCwyLCJcXHZEYmxDYXQiXSxbMiwwLCJcXHZEYmxDYXQiXSxbMiwyLCJcXHZEYmxDYXQiXSxbMCwyLCJcXE1vZCgtKSJdLFsxLDMsIlxcTW9kKC0pIiwyXSxbMCwxLCJcXG9wX0giLDJdLFsyLDMsIlxcb3BfSCJdLFsyLDEsInIiLDAseyJzaG9ydGVuIjp7InNvdXJjZSI6MjAsInRhcmdldCI6MjB9LCJsZXZlbCI6Mn1dXQ==
	\[\begin{tikzcd}[ampersand replacement=\&]
		\vDblCat \&\& \vDblCat \\
		\\
		\vDblCat \&\& \vDblCat
		\arrow["{\Mod(-)}", from=1-1, to=1-3]
		\arrow["{\op_H}"', from=1-1, to=3-1]
		\arrow["r", between={0.2}{0.8}, Rightarrow, nfold, from=1-3, to=3-1]
		\arrow["{\op_H}", from=1-3, to=3-3]
		\arrow["{\Mod(-)}"', from=3-1, to=3-3]
	\end{tikzcd}\cdpunct[8pt]{.}\]
\end{theorem}

\begin{proof} The components $r_{\bV} \colon \op_H(\Mod(\bV)) \rightarrow \Mod(\op_{H}(\bV))$ are defined by 
	\begin{description}
		\item[Objects] 
		% https://q.uiver.app/#q=WzAsMTgsWzEsMCwiTSBcXGNvbG9uZXFxIChcXGJ1bGxldCxNLFxcbXUsZSkiXSxbNSwwLCJNXnIgXFxjb2xvbmVxcSAoXFxidWxsZXQsTV5yLFxcbXUsZSkiXSxbMiwyLCJcXGJ1bGxldCJdLFsyLDEsIlxcYnVsbGV0Il0sWzEsMSwiXFxidWxsZXQiXSxbMCwxLCJcXGJ1bGxldCJdLFswLDIsIlxcYnVsbGV0Il0sWzQsMSwiXFxidWxsZXQiXSxbNCwyLCJcXGJ1bGxldCJdLFs1LDEsIlxcYnVsbGV0Il0sWzYsMSwiXFxidWxsZXQiXSxbNiwyLCJcXGJ1bGxldCJdLFsyLDQsIlxcYnVsbGV0Il0sWzAsNCwiXFxidWxsZXQiXSxbMSwzLCJcXGJ1bGxldCJdLFs0LDQsIlxcYnVsbGV0Il0sWzYsNCwiXFxidWxsZXQiXSxbNSwzLCJcXGJ1bGxldCJdLFszLDQsIk0iLDIseyJzdHlsZSI6eyJib2R5Ijp7Im5hbWUiOiJiYXJyZWQifX19XSxbNCw1LCJNIiwyLHsic3R5bGUiOnsiYm9keSI6eyJuYW1lIjoiYmFycmVkIn19fV0sWzIsMywiIiwwLHsic3R5bGUiOnsiaGVhZCI6eyJuYW1lIjoibm9uZSJ9fX1dLFsyLDYsIk0iLDAseyJzdHlsZSI6eyJib2R5Ijp7Im5hbWUiOiJiYXJyZWQifX19XSxbNiw1LCIiLDIseyJzdHlsZSI6eyJoZWFkIjp7Im5hbWUiOiJub25lIn19fV0sWzcsOSwiTV5yIiwwLHsic3R5bGUiOnsiYm9keSI6eyJuYW1lIjoiYmFycmVkIn19fV0sWzksMTAsIk1eciIsMCx7InN0eWxlIjp7ImJvZHkiOnsibmFtZSI6ImJhcnJlZCJ9fX1dLFsxMCwxMSwiIiwwLHsic3R5bGUiOnsiaGVhZCI6eyJuYW1lIjoibm9uZSJ9fX1dLFs4LDExLCJNXnIiLDIseyJzdHlsZSI6eyJib2R5Ijp7Im5hbWUiOiJiYXJyZWQifX19XSxbNyw4LCIiLDIseyJzdHlsZSI6eyJoZWFkIjp7Im5hbWUiOiJub25lIn19fV0sWzEyLDEzLCJNIiwwLHsic3R5bGUiOnsiYm9keSI6eyJuYW1lIjoiYmFycmVkIn19fV0sWzEyLDE0LCIiLDIseyJzdHlsZSI6eyJoZWFkIjp7Im5hbWUiOiJub25lIn19fV0sWzEzLDE0LCIiLDAseyJzdHlsZSI6eyJoZWFkIjp7Im5hbWUiOiJub25lIn19fV0sWzE3LDE1LCIiLDIseyJzdHlsZSI6eyJoZWFkIjp7Im5hbWUiOiJub25lIn19fV0sWzE3LDE2LCIiLDAseyJzdHlsZSI6eyJoZWFkIjp7Im5hbWUiOiJub25lIn19fV0sWzE1LDE2LCJNXnIiLDIseyJzdHlsZSI6eyJib2R5Ijp7Im5hbWUiOiJiYXJyZWQifX19XSxbMCwxLCIiLDAseyJzaG9ydGVuIjp7InNvdXJjZSI6MjAsInRhcmdldCI6MjB9LCJzdHlsZSI6eyJ0YWlsIjp7Im5hbWUiOiJtYXBzIHRvIn19fV0sWzQsMjEsIlxcbXUiLDIseyJzaG9ydGVuIjp7InRhcmdldCI6MjB9fV0sWzksMjYsIlxcbXUiLDAseyJzaG9ydGVuIjp7InRhcmdldCI6MjB9fV0sWzE0LDI4LCJlIiwyLHsic2hvcnRlbiI6eyJ0YXJnZXQiOjIwfX1dLFsxNywzMywiZSIsMCx7InNob3J0ZW4iOnsic291cmNlIjoyMCwidGFyZ2V0IjoyMH19XV0=
		\[\begin{tikzcd}[ampersand replacement=\&]
			\& {M \coloneqq (\bullet,M,\mu,e)} \&\&\&\& {M^r \coloneqq (\bullet,M^r,\mu,e)} \& \\
			\bullet \& \bullet \& \bullet \&\& \bullet \& \bullet \& \bullet \\
			\bullet \&\& \bullet \&\& \bullet \&\& \bullet \\
			\& \bullet \&\&\&\& \bullet \\
			\bullet \&\& \bullet \&\& \bullet \&\& \bullet
			\arrow[between={0.2}{0.8}, maps to, from=1-2, to=1-6]
			\arrow["M"'{inner sep=.8ex}, "\shortmid"{marking}, from=2-2, to=2-1]
			\arrow["M"'{inner sep=.8ex}, "\shortmid"{marking}, from=2-3, to=2-2]
			\arrow["{M^r}"{inner sep=.8ex}, "\shortmid"{marking}, from=2-5, to=2-6]
			\arrow[no head, from=2-5, to=3-5]
			\arrow["{M^r}"{inner sep=.8ex}, "\shortmid"{marking}, from=2-6, to=2-7]
			\arrow[no head, from=2-7, to=3-7]
			\arrow[no head, from=3-1, to=2-1]
			\arrow[no head, from=3-3, to=2-3]
			\arrow[""{name=0, anchor=center, inner sep=0}, "M"{inner sep=.8ex}, "\shortmid"{marking}, from=3-3, to=3-1]
			\arrow[""{name=1, anchor=center, inner sep=0}, "{M^r}"'{inner sep=.8ex}, "\shortmid"{marking}, from=3-5, to=3-7]
			\arrow[no head, from=4-6, to=5-5]
			\arrow[no head, from=4-6, to=5-7]
			\arrow[no head, from=5-1, to=4-2]
			\arrow[no head, from=5-3, to=4-2]
			\arrow[""{name=2, anchor=center, inner sep=0}, "M"{inner sep=.8ex}, "\shortmid"{marking}, from=5-3, to=5-1]
			\arrow[""{name=3, anchor=center, inner sep=0}, "{M^r}"'{inner sep=.8ex}, "\shortmid"{marking}, from=5-5, to=5-7]
			\arrow["\mu"', between={0}{0.8}, Rightarrow, nfold, from=2-2, to=0]
			\arrow["\mu", between={0}{0.8}, Rightarrow, nfold, from=2-6, to=1]
			\arrow["e"', between={0}{0.8}, Rightarrow, nfold, from=4-2, to=2]
			\arrow["e", between={0.2}{0.8}, Rightarrow, nfold, from=4-6, to=3]
		\end{tikzcd}\]
		\item[V $1$-cell] 
		% https://q.uiver.app/#q=WzAsOCxbMiwwLCJcXGJ1bGxldCJdLFswLDAsIlxcYnVsbGV0Il0sWzIsMiwiXFxzdGFyIl0sWzAsMiwiXFxzdGFyIl0sWzUsMCwiXFxidWxsZXQiXSxbNywwLCJcXGJ1bGxldCJdLFs1LDIsIlxcc3RhciJdLFs3LDIsIlxcc3RhciJdLFswLDEsIk0iLDIseyJzdHlsZSI6eyJib2R5Ijp7Im5hbWUiOiJiYXJyZWQifX19XSxbMCwyLCJmX3swfSJdLFsxLDMsImZfezB9IiwyXSxbMiwzLCJOIiwwLHsic3R5bGUiOnsiYm9keSI6eyJuYW1lIjoiYmFycmVkIn19fV0sWzYsNywiTl5yIiwyLHsic3R5bGUiOnsiYm9keSI6eyJuYW1lIjoiYmFycmVkIn19fV0sWzQsNSwiTV5yIiwwLHsic3R5bGUiOnsiYm9keSI6eyJuYW1lIjoiYmFycmVkIn19fV0sWzQsNiwiZl8wIiwyXSxbNSw3LCJmXzAiXSxbOCwxMSwiZiIsMCx7InNob3J0ZW4iOnsic291cmNlIjoyMCwidGFyZ2V0IjoyMH19XSxbMTMsMTIsImYiLDAseyJzaG9ydGVuIjp7InNvdXJjZSI6MjAsInRhcmdldCI6MjB9fV0sWzksMTQsIiIsMix7InNob3J0ZW4iOnsic291cmNlIjoyMCwidGFyZ2V0IjoyMH0sImxldmVsIjoxLCJzdHlsZSI6eyJ0YWlsIjp7Im5hbWUiOiJtYXBzIHRvIn19fV1d
		\[\begin{tikzcd}[ampersand replacement=\&]
			\bullet \&\& \bullet \&\&\& \bullet \&\& \bullet \\
			\\
			\star \&\& \star \&\&\& \star \&\& \star
			\arrow["{f_{0}}"', from=1-1, to=3-1]
			\arrow[""{name=0, anchor=center, inner sep=0}, "M"'{inner sep=.8ex}, "\shortmid"{marking}, from=1-3, to=1-1]
			\arrow[""{name=1, anchor=center, inner sep=0}, "{f_{0}}", from=1-3, to=3-3]
			\arrow[""{name=2, anchor=center, inner sep=0}, "{M^r}"{inner sep=.8ex}, "\shortmid"{marking}, from=1-6, to=1-8]
			\arrow[""{name=3, anchor=center, inner sep=0}, "{f_0}"', from=1-6, to=3-6]
			\arrow["{f_0}", from=1-8, to=3-8]
			\arrow[""{name=4, anchor=center, inner sep=0}, "N"{inner sep=.8ex}, "\shortmid"{marking}, from=3-3, to=3-1]
			\arrow[""{name=5, anchor=center, inner sep=0}, "{N^r}"'{inner sep=.8ex}, "\shortmid"{marking}, from=3-6, to=3-8]
			\arrow["f", between={0.2}{0.8}, Rightarrow, nfold, from=0, to=4]
			\arrow[between={0.2}{0.8}, maps to, from=1, to=3]
			\arrow["f", between={0.2}{0.8}, Rightarrow, nfold, from=2, to=5]
		\end{tikzcd}\]
		\item[H $1$-cell]
		% https://q.uiver.app/#q=WzAsMjgsWzAsMSwiXFxzdGFyIl0sWzQsMywiXFxidWxsZXQiXSxbNCwyLCJcXGJ1bGxldCJdLFsyLDIsIlxcYnVsbGV0Il0sWzAsMiwiXFxzdGFyIl0sWzAsMywiXFxzdGFyIl0sWzQsNCwiXFxidWxsZXQiXSxbNCw1LCJcXGJ1bGxldCJdLFsyLDQsIlxcc3RhciJdLFswLDQsIlxcc3RhciJdLFswLDUsIlxcc3RhciJdLFs0LDEsIlxcYnVsbGV0Il0sWzYsMSwiXFxzdGFyIl0sWzEwLDEsIlxcYnVsbGV0Il0sWzEwLDIsIlxcYnVsbGV0Il0sWzgsMiwiXFxidWxsZXQiXSxbMTAsMywiXFxidWxsZXQiXSxbOCw0LCJcXHN0YXIiXSxbNiw0LCJcXHN0YXIiXSxbNiw1LCJcXHN0YXIiXSxbMTAsNSwiXFxidWxsZXQiXSxbMTAsNCwiXFxidWxsZXQiXSxbNiwzLCJcXHN0YXIiXSxbNiwyLCJcXHN0YXIiXSxbNCwwLCIoTSxcXGJ1bGxldCkiXSxbMCwwLCIoTixcXHN0YXIpIl0sWzYsMCwiKE5ecixcXHN0YXIpIl0sWzEwLDAsIihNXnIsXFxidWxsZXQpIl0sWzIsMywiTSIsMix7InN0eWxlIjp7ImJvZHkiOnsibmFtZSI6ImJhcnJlZCJ9fX1dLFszLDQsIlgiLDIseyJzdHlsZSI6eyJib2R5Ijp7Im5hbWUiOiJiYXJyZWQifX19XSxbMSwyLCIiLDAseyJzdHlsZSI6eyJoZWFkIjp7Im5hbWUiOiJub25lIn19fV0sWzEsNSwiWCIsMCx7InN0eWxlIjp7ImJvZHkiOnsibmFtZSI6ImJhcnJlZCJ9fX1dLFs1LDQsIiIsMix7InN0eWxlIjp7ImhlYWQiOnsibmFtZSI6Im5vbmUifX19XSxbNiw4LCJYIiwyLHsic3R5bGUiOnsiYm9keSI6eyJuYW1lIjoiYmFycmVkIn19fV0sWzgsOSwiTiIsMix7InN0eWxlIjp7ImJvZHkiOnsibmFtZSI6ImJhcnJlZCJ9fX1dLFs5LDEwLCIiLDAseyJzdHlsZSI6eyJoZWFkIjp7Im5hbWUiOiJub25lIn19fV0sWzcsMTAsIlgiLDAseyJzdHlsZSI6eyJib2R5Ijp7Im5hbWUiOiJiYXJyZWQifX19XSxbNiw3LCIiLDIseyJzdHlsZSI6eyJoZWFkIjp7Im5hbWUiOiJub25lIn19fV0sWzExLDAsIlgiLDIseyJzdHlsZSI6eyJib2R5Ijp7Im5hbWUiOiJiYXJyZWQifX19XSxbMTIsMTMsIlheciIsMCx7InN0eWxlIjp7ImJvZHkiOnsibmFtZSI6ImJhcnJlZCJ9fX1dLFsyMiwxNiwiWF5yIiwyLHsic3R5bGUiOnsiYm9keSI6eyJuYW1lIjoiYmFycmVkIn19fV0sWzE1LDE0LCJNXnIiLDAseyJzdHlsZSI6eyJib2R5Ijp7Im5hbWUiOiJiYXJyZWQifX19XSxbMTcsMjEsIlheciIsMCx7InN0eWxlIjp7ImJvZHkiOnsibmFtZSI6ImJhcnJlZCJ9fX1dLFsxOCwxNywiTl5yIiwwLHsic3R5bGUiOnsiYm9keSI6eyJuYW1lIjoiYmFycmVkIn19fV0sWzE5LDIwLCJYXnIiLDIseyJzdHlsZSI6eyJib2R5Ijp7Im5hbWUiOiJiYXJyZWQifX19XSxbMjMsMTUsIlheciIsMCx7InN0eWxlIjp7ImJvZHkiOnsibmFtZSI6ImJhcnJlZCJ9fX1dLFsxNCwxNiwiIiwxLHsic3R5bGUiOnsiaGVhZCI6eyJuYW1lIjoibm9uZSJ9fX1dLFsyMywyMiwiIiwxLHsic3R5bGUiOnsiaGVhZCI6eyJuYW1lIjoibm9uZSJ9fX1dLFsxOCwxOSwiIiwxLHsic3R5bGUiOnsiaGVhZCI6eyJuYW1lIjoibm9uZSJ9fX1dLFsyMSwyMCwiIiwxLHsic3R5bGUiOnsiaGVhZCI6eyJuYW1lIjoibm9uZSJ9fX1dLFsyNiwyNywiWF57cn1cXGNvbG9uZXFxKFhecixcXHJobyxcXGxhbWJkYSkiLDAseyJzdHlsZSI6eyJib2R5Ijp7Im5hbWUiOiJiYXJyZWQifX19XSxbMjQsMjUsIlggXFxjb2xvbmVxcShYLFxcbGFtYmRhLFxccmhvKSIsMix7InN0eWxlIjp7ImJvZHkiOnsibmFtZSI6ImJhcnJlZCJ9fX1dLFsyNCwyNiwiIiwwLHsic2hvcnRlbiI6eyJzb3VyY2UiOjIwLCJ0YXJnZXQiOjIwfSwic3R5bGUiOnsidGFpbCI6eyJuYW1lIjoibWFwcyB0byJ9fX1dLFszLDMxLCJcXHJobyIsMix7InNob3J0ZW4iOnsidGFyZ2V0IjoyMH19XSxbOCwzNiwiXFxsYW1iZGEiLDIseyJzaG9ydGVuIjp7InRhcmdldCI6MjB9fV0sWzE1LDQwLCJcXHJobyIsMix7InNob3J0ZW4iOnsidGFyZ2V0IjoyMH19XSxbMTcsNDQsIlxcbGFtYmRhIiwyLHsic2hvcnRlbiI6eyJ0YXJnZXQiOjIwfX1dXQ==
		\[\begin{tikzcd}[ampersand replacement=\&,column sep=small]
			{(N,\star)} \&\&\&\& {(M,\bullet)} \&\& {(N^r,\star)} \&\&\&\& {(M^r,\bullet)} \\
			\star \&\&\&\& \bullet \&\& \star \&\&\&\& \bullet \\
			\star \&\& \bullet \&\& \bullet \&\& \star \&\& \bullet \&\& \bullet \\
			\star \&\&\&\& \bullet \&\& \star \&\&\&\& \bullet \\
			\star \&\& \star \&\& \bullet \&\& \star \&\& \star \&\& \bullet \\
			\star \&\&\&\& \bullet \&\& \star \&\&\&\& \bullet
			\arrow["{X \coloneqq(X,\lambda,\rho)}"'{inner sep=.8ex}, "\shortmid"{marking}, from=1-5, to=1-1]
			\arrow[between={0.2}{0.8}, maps to, from=1-5, to=1-7]
			\arrow["{X^{r}\coloneqq(X^r,\rho,\lambda)}"{inner sep=.8ex}, "\shortmid"{marking}, from=1-7, to=1-11]
			\arrow["X"'{inner sep=.8ex}, "\shortmid"{marking}, from=2-5, to=2-1]
			\arrow["{X^r}"{inner sep=.8ex}, "\shortmid"{marking}, from=2-7, to=2-11]
			\arrow["X"'{inner sep=.8ex}, "\shortmid"{marking}, from=3-3, to=3-1]
			\arrow["M"'{inner sep=.8ex}, "\shortmid"{marking}, from=3-5, to=3-3]
			\arrow["{X^r}"{inner sep=.8ex}, "\shortmid"{marking}, from=3-7, to=3-9]
			\arrow[no head, from=3-7, to=4-7]
			\arrow["{M^r}"{inner sep=.8ex}, "\shortmid"{marking}, from=3-9, to=3-11]
			\arrow[no head, from=3-11, to=4-11]
			\arrow[no head, from=4-1, to=3-1]
			\arrow[no head, from=4-5, to=3-5]
			\arrow[""{name=0, anchor=center, inner sep=0}, "X"{inner sep=.8ex}, "\shortmid"{marking}, from=4-5, to=4-1]
			\arrow[""{name=1, anchor=center, inner sep=0}, "{X^r}"'{inner sep=.8ex}, "\shortmid"{marking}, from=4-7, to=4-11]
			\arrow[no head, from=5-1, to=6-1]
			\arrow["N"'{inner sep=.8ex}, "\shortmid"{marking}, from=5-3, to=5-1]
			\arrow["X"'{inner sep=.8ex}, "\shortmid"{marking}, from=5-5, to=5-3]
			\arrow[no head, from=5-5, to=6-5]
			\arrow["{N^r}"{inner sep=.8ex}, "\shortmid"{marking}, from=5-7, to=5-9]
			\arrow[no head, from=5-7, to=6-7]
			\arrow["{X^r}"{inner sep=.8ex}, "\shortmid"{marking}, from=5-9, to=5-11]
			\arrow[no head, from=5-11, to=6-11]
			\arrow[""{name=2, anchor=center, inner sep=0}, "X"{inner sep=.8ex}, "\shortmid"{marking}, from=6-5, to=6-1]
			\arrow[""{name=3, anchor=center, inner sep=0}, "{X^r}"'{inner sep=.8ex}, "\shortmid"{marking}, from=6-7, to=6-11]
			\arrow["\rho"', between={0}{0.8}, Rightarrow, nfold, from=3-3, to=0]
			\arrow["\rho"', between={0}{0.8}, Rightarrow, nfold, from=3-9, to=1]
			\arrow["\lambda"', between={0}{0.8}, Rightarrow, nfold, from=5-3, to=2]
			\arrow["\lambda"', between={0}{0.8}, Rightarrow, nfold, from=5-9, to=3]
		\end{tikzcd}\]
		\item[Sq multicell]
		% https://q.uiver.app/#q=WzAsMTAsWzQsMiwiXFxidWxsZXQiXSxbMCwyLCJcXHN0YXIiXSxbNCwwLCJcXGJ1bGxldF8xIl0sWzAsMCwiXFxidWxsZXRfbiJdLFsyLDAsIlxcY2RvdHMiXSxbNiwwLCJcXGJ1bGxldF9uIl0sWzgsMCwiXFxjZG90cyJdLFsxMCwwLCJcXGJ1bGxldF8xIl0sWzYsMiwiXFxzdGFyIl0sWzEwLDIsIlxcYnVsbGV0Il0sWzIsNCwiWF8xIiwyLHsic3R5bGUiOnsiYm9keSI6eyJuYW1lIjoiYmFycmVkIn19fV0sWzQsMywiWF9uIiwyLHsic3R5bGUiOnsiYm9keSI6eyJuYW1lIjoiYmFycmVkIn19fV0sWzAsMSwiWSIsMCx7InN0eWxlIjp7ImJvZHkiOnsibmFtZSI6ImJhcnJlZCJ9fX1dLFsyLDAsImYiXSxbMywxLCJnIiwyXSxbNSw4LCJnIiwyXSxbNyw5LCJmIl0sWzgsOSwiWV5yIiwyLHsic3R5bGUiOnsiYm9keSI6eyJuYW1lIjoiYmFycmVkIn19fV0sWzUsNiwiWF9uXnIiLDAseyJzdHlsZSI6eyJib2R5Ijp7Im5hbWUiOiJiYXJyZWQifX19XSxbNiw3LCJYXzFeciIsMCx7InN0eWxlIjp7ImJvZHkiOnsibmFtZSI6ImJhcnJlZCJ9fX1dLFs0LDEyLCJcXGFscGhhIiwwLHsic2hvcnRlbiI6eyJzb3VyY2UiOjEwLCJ0YXJnZXQiOjIwfX1dLFs2LDE3LCJcXGFscGhhIiwwLHsic2hvcnRlbiI6eyJzb3VyY2UiOjEwLCJ0YXJnZXQiOjIwfX1dLFsxMywxNSwiIiwyLHsic2hvcnRlbiI6eyJzb3VyY2UiOjMwLCJ0YXJnZXQiOjMwfSwibGV2ZWwiOjEsInN0eWxlIjp7InRhaWwiOnsibmFtZSI6Im1hcHMgdG8ifX19XV0=
		\[\begin{tikzcd}[ampersand replacement=\&]
			{\bullet_n} \&\& \cdots \&\& {\bullet_1} \&\& {\bullet_n} \&\& \cdots \&\& {\bullet_1} \\
			\\
			\star \&\&\&\& \bullet \&\& \star \&\&\&\& \bullet
			\arrow["g"', from=1-1, to=3-1]
			\arrow["{X_n}"'{inner sep=.8ex}, "\shortmid"{marking}, from=1-3, to=1-1]
			\arrow["{X_1}"'{inner sep=.8ex}, "\shortmid"{marking}, from=1-5, to=1-3]
			\arrow[""{name=0, anchor=center, inner sep=0}, "f", from=1-5, to=3-5]
			\arrow["{X_n^r}"{inner sep=.8ex}, "\shortmid"{marking}, from=1-7, to=1-9]
			\arrow[""{name=1, anchor=center, inner sep=0}, "g"', from=1-7, to=3-7]
			\arrow["{X_1^r}"{inner sep=.8ex}, "\shortmid"{marking}, from=1-9, to=1-11]
			\arrow["f", from=1-11, to=3-11]
			\arrow[""{name=2, anchor=center, inner sep=0}, "Y"{inner sep=.8ex}, "\shortmid"{marking}, from=3-5, to=3-1]
			\arrow[""{name=3, anchor=center, inner sep=0}, "{Y^r}"'{inner sep=.8ex}, "\shortmid"{marking}, from=3-7, to=3-11]
			\arrow["\alpha", between={0.1}{0.8}, Rightarrow, nfold, from=1-3, to=2]
			\arrow[between={0.3}{0.7}, maps to, from=0, to=1]
			\arrow["\alpha", between={0.1}{0.8}, Rightarrow, nfold, from=1-9, to=3]
		\end{tikzcd}\]
	\end{description}
\end{proof}

We shall now specialize back to the case of monoidal categories. In particular, we shall restrict $\Mod(-)$ and $\CoMod(-)$ to monoidal categories.  In what follows, we write our diagrams in strict monoidal categories $\st \cV$ and interpret the $2$-cells using the canonical evaluation map.

\begin{comment}
	To this end, it is useful to express the square multicells of $\Mod(\cV)$ explicitly.
	
	\begin{definition}[Balancing law]
		Let $Y$ be a $(M,N)$-module in $\cV$. Let $A$ be a $(M,P)$-module and let $B$ be a $(P,N)$-module similarly. A morphism $f \colon A \otimes B \rightarrow Y$ satisfy \emph{balancing law} if it coequalizes the two actions % https://q.uiver.app/#q=WzAsMixbMiwwLCJBIFxcb3RpbWVzIEIiXSxbMCwwLCJBIFxcb3RpbWVzIFAgXFxvdGltZXMgQiJdLFsxLDAsIjEgXFxvdGltZXMgXFxsYW1iZGFfe0J9IiwwLHsib2Zmc2V0IjotMn1dLFsxLDAsIlxccmhvX3tBfSBcXG90aW1lcyAxIiwyLHsib2Zmc2V0IjoyfV1d
		\begin{tikzcd}[ampersand replacement=\&,cramped]
			{A \otimes P \otimes B} \&\& {A \otimes B}
			\arrow["{1 \otimes \lambda_{B}}", shift left=2, from=1-1, to=1-3]
			\arrow["{\rho_{A} \otimes 1}"', shift right=2, from=1-1, to=1-3]
		\end{tikzcd}.
	\end{definition}
	
	Of course, if $f$ universal then it simply describes tensor product map $A \otimes B \rightarrow A \otimes_{P} B$. However, as we have already stated, we do not wish to assume completeness properties. In this context, we still need a notion of iterated tensor product of modules. To this end, we note that the two parallel arrows can be captured as a functor $\mathrm{bl}_{A,P,B} \colon \bDelta^{\op}_{\leq 1} \rightarrow \cV$. Thereafter, the morphism $f$ is merely a cocone under $\mathrm{bl}_{A,P,B}$. This suggests the following more general definition.
	
	\begin{definition}[Multi-balancing law]
		Consider a string of $n$-modules $A=(A_1, \cdots ,A_n)$ such that each $A_i$ and $A_{i+1}$ share the same intermediary monoid. Let $\mathrm{bl}_{A}\colon \left(\bDelta_{\leq 1}^{\op}\right)^{n} \rightarrow \cV$ be the canonical functor, which for each copy of $\bDelta_{\leq 1}^{\op}$ is defined on the atomic arrows as the parallel pair of the corresponding module actions. We say that a morphism $f : A_1 \otimes A_2 \otimes \cdots \otimes A_n \rightarrow B$ satisfy \emph{multi-balancing law} if it extends to a cocone under the given functor. 
	\end{definition}
	
	Finally, suppose that the string of modules $A = (A_1, \cdots, A_n)$ is such that $A_1$ is left $M$-module and $A_n$ is right $N$-module. If there are morphisms of monoid $f_{M,M'}\colon :M \rightarrow M'$, $f_{N,N'} \colon N \rightarrow N'$ and $B$ is a $(M',N')$-module then we also need to incorporate the action of the monoid morphisms into our diagram. Let $f\colon A_1 \otimes \cdots \otimes A_n \rightarrow B$ be a multi-balanced morphism. We require that $f$ satisfy the following two additional identites :
	% https://q.uiver.app/#q=WzAsOCxbMCwwLCJNIFxcb3RpbWVzIEEiXSxbMCwyLCJBIl0sWzIsMiwiQiJdLFsyLDAsIk0nIFxcb3RpbWVzIEIiXSxbNCwwLCJBIFxcb3RpbWVzIE4iXSxbNiwwLCJCIFxcb3RpbWVzIE4nIl0sWzQsMiwiQSJdLFs2LDIsIkIiXSxbMCwxLCJcXGxhbWJkYV97QV8xfSIsMl0sWzEsMiwiZiIsMl0sWzAsMywiZl97TSxNJ30gXFxvdGltZXMgZiJdLFszLDIsIlxcbGFtYmRhX3tZfSJdLFs0LDYsIlxccmhvX3tBX3tufX0iLDJdLFs2LDcsImYiLDJdLFs0LDUsImYgXFxvdGltZXMgZl97TixOJ30iXSxbNSw3LCJcXHJob197WX0iXV0=
	\[\begin{tikzcd}[ampersand replacement=\&,cramped]
		{M \otimes A} \&\& {M' \otimes B} \&\& {A \otimes N} \&\& {B \otimes N'} \\
		\\
		A \&\& B \&\& A \&\& B
		\arrow["{f_{M,M'} \otimes f}", from=1-1, to=1-3]
		\arrow["{\lambda_{A_1}}"', from=1-1, to=3-1]
		\arrow["{\lambda_{Y}}", from=1-3, to=3-3]
		\arrow["{f \otimes f_{N,N'}}", from=1-5, to=1-7]
		\arrow["{\rho_{A_{n}}}"', from=1-5, to=3-5]
		\arrow["{\rho_{Y}}", from=1-7, to=3-7]
		\arrow["f"', from=3-1, to=3-3]
		\arrow["f"', from=3-5, to=3-7]
	\end{tikzcd}\cdpunct[8pt]{.}\]
	
	For the case of $\vcoDblCat$, the dual multi-balancing laws are expressed in terms of equalizers.
\end{comment}

\begin{proposition}
	The $(\op,\op^t)$ morphism of inverses on monoidal categories extends to virtual double categories through the following strict $2$-transformation cells:
	% https://q.uiver.app/#q=WzAsOCxbMCwwLCJcXG9wXzIoXFxNb25DYXRfYykiXSxbMCwyLCJcXE1vbkNhdF9sIl0sWzMsMiwiXFx2RGJsQ2F0Il0sWzMsMCwiXFxvcF8yKFxcdmNvRGJsQ2F0KSJdLFswLDMsIlxcb3BfMihcXE1vbkNhdF9sKSJdLFswLDUsIlxcTW9uQ2F0X2MiXSxbMywzLCJcXG9wXzIoXFx2RGJsQ2F0KSJdLFszLDUsIlxcdmNvRGJsQ2F0Il0sWzAsMSwiXFxvcF8xIiwyXSxbMywyLCJcXG9wX1YiXSxbMSwyLCJcXE1vZCgtKSIsMl0sWzAsMywiXFxvcF8yKFxcQ29Nb2QoLSkpIl0sWzMsMSwiXFxvcCIsMCx7InNob3J0ZW4iOnsic291cmNlIjoyMCwidGFyZ2V0IjoyMH0sImxldmVsIjoyfV0sWzQsNSwiXFxvcF8xXnt0fSIsMl0sWzQsNiwiXFxvcF8yKFxcTW9kKC0pKSJdLFs2LDcsIlxcb3Bfe1Z9Xnt0fSJdLFs1LDcsIlxcQ29Nb2QoLSkiLDJdLFs2LDUsIlxcb3Bee3R9IiwwLHsic2hvcnRlbiI6eyJzb3VyY2UiOjIwLCJ0YXJnZXQiOjIwfSwibGV2ZWwiOjJ9XV0=
	\[\begin{tikzcd}[ampersand replacement=\&]
		{\op_2(\MonCat_c)} \&\&\& {\op_2(\vcoDblCat)} \\
		\\
		{\MonCat_l} \&\&\& \vDblCat \\
		{\op_2(\MonCat_l)} \&\&\& {\op_2(\vDblCat)} \\
		\\
		{\MonCat_c} \&\&\& \vcoDblCat
		\arrow["{\op_2(\CoMod(-))}", from=1-1, to=1-4]
		\arrow["{\op_1}"', from=1-1, to=3-1]
		\arrow["\op", between={0.2}{0.8}, Rightarrow, nfold, from=1-4, to=3-1]
		\arrow["{\op_V}", from=1-4, to=3-4]
		\arrow["{\Mod(-)}"', from=3-1, to=3-4]
		\arrow["{\op_2(\Mod(-))}", from=4-1, to=4-4]
		\arrow["{\op_1^{t}}"', from=4-1, to=6-1]
		\arrow["{\op^{t}}", between={0.2}{0.8}, Rightarrow, nfold, from=4-4, to=6-1]
		\arrow["{\op_{V}^{t}}", from=4-4, to=6-4]
		\arrow["{\CoMod(-)}"', from=6-1, to=6-4]
	\end{tikzcd}\cdpunct[8pt]{.}\]
\end{proposition}

\begin{proof}
	We primarily need to compute the components $\op_{\cV}$ and $\op^t_{\cV}$. For the component $\op_{\cV}$, this means that we need to obtain a virtual double functor $\op_{\cV}\colon \op_{V}(\CoMod(\cV)) \rightarrow \Mod(\op_{1}(\cV))$. This is given by the following data:
	\begin{description}
		\item[Objects] 
		% https://q.uiver.app/#q=WzAsMixbMCwwLCJDIFxcY29sb25lcXEgKEMsXFxtdSxlKSJdLFsyLDAsIkNee1xcb3B9IFxcY29sb25lcXEgKEMsIFxcbXVee1xcb3B9LCBlXntcXG9wfSkiXSxbMCwxLCIiLDAseyJzdHlsZSI6eyJ0YWlsIjp7Im5hbWUiOiJtYXBzIHRvIn19fV1d
		\[\begin{tikzcd}[ampersand replacement=\&]
			{C \coloneqq (C,\mu,e)} \&\& {C^{\op} \coloneqq (C, \mu^{\op}, e^{\op})}
			\arrow[maps to, from=1-1, to=1-3]
		\end{tikzcd}\]
		\item[V $1$-cell] 
		% https://q.uiver.app/#q=WzAsNCxbMCwyLCJDIl0sWzAsMCwiQyciXSxbMywwLCIoQycpXntcXG9wfSJdLFszLDIsIkNee1xcb3B9Il0sWzAsMV0sWzIsM10sWzQsNSwiIiwwLHsic2hvcnRlbiI6eyJzb3VyY2UiOjIwLCJ0YXJnZXQiOjIwfSwibGV2ZWwiOjEsInN0eWxlIjp7InRhaWwiOnsibmFtZSI6Im1hcHMgdG8ifX19XV0=
		\[\begin{tikzcd}[ampersand replacement=\&]
			{C'} \&\&\& {(C')^{\op}} \\
			\\
			C \&\&\& {C^{\op}}
			\arrow[""{name=0, anchor=center, inner sep=0}, from=1-4, to=3-4]
			\arrow[""{name=1, anchor=center, inner sep=0}, from=3-1, to=1-1]
			\arrow[between={0.2}{0.8}, maps to, from=1, to=0]
		\end{tikzcd}\]
		\item[H $1$-cell] 
		% https://q.uiver.app/#q=WzAsNCxbMCwwLCJDXzEiXSxbMywwLCJDXzIiXSxbNiwwLCJDXzFee1xcb3B9Il0sWzksMCwiQ18yXntcXG9wfSJdLFswLDEsIlggXFxjb2xvbmVxcShYLFxcbGFtYmRhLFxccmhvKSIsMCx7InN0eWxlIjp7ImJvZHkiOnsibmFtZSI6ImJhcnJlZCJ9fX1dLFsyLDMsIlhee1xcb3B9IFxcY29sb25lcXEgKFggLFxcbGFtYmRhXntcXG9wfSxcXHJob157XFxvcH0pIiwwLHsic3R5bGUiOnsiYm9keSI6eyJuYW1lIjoiYmFycmVkIn19fV0sWzQsNSwiIiwwLHsic2hvcnRlbiI6eyJzb3VyY2UiOjQwLCJ0YXJnZXQiOjQwfSwibGV2ZWwiOjEsInN0eWxlIjp7InRhaWwiOnsibmFtZSI6Im1hcHMgdG8ifX19XV0=
		\[\begin{tikzcd}[ampersand replacement=\&]
			{C_1} \&\&\& {C_2} \&\&\& {C_1^{\op}} \&\&\& {C_2^{\op}}
			\arrow[""{name=0, anchor=center, inner sep=0}, "{X \coloneqq(X,\lambda,\rho)}"{inner sep=.8ex}, "\shortmid"{marking}, from=1-1, to=1-4]
			\arrow[""{name=1, anchor=center, inner sep=0}, "{X^{\op} \coloneqq (X ,\lambda^{\op},\rho^{\op})}"{inner sep=.8ex}, "\shortmid"{marking}, from=1-7, to=1-10]
			\arrow[between={0.4}{0.6}, maps to, from=0, to=1]
		\end{tikzcd}\]
		\item[Sq multicell] 
		% https://q.uiver.app/#q=WzAsMTIsWzAsMiwiQ18xIl0sWzMsMiwiQ18yIl0sWzAsMCwiQ18xJyJdLFszLDAsIkNfMiciXSxbMSwwLCJcXGNkb3RzIl0sWzIsMCwiXFxjZG90cyJdLFs1LDAsIihDXzEnKV57XFxvcH0iXSxbNSwyLCJDXzFee1xcb3B9Il0sWzgsMiwiQ18yXntcXG9wfSJdLFs2LDAsIlxcY2RvdHMiXSxbNywwLCJcXGNkb3RzIl0sWzgsMCwiKENfMicpXntcXG9wfSJdLFswLDEsIlgiLDIseyJzdHlsZSI6eyJib2R5Ijp7Im5hbWUiOiJiYXJyZWQifX19XSxbMCwyXSxbMSwzXSxbMiw0LCJZIiwwLHsic3R5bGUiOnsiYm9keSI6eyJuYW1lIjoiYmFycmVkIn19fV0sWzUsMywiWiIsMCx7InN0eWxlIjp7ImJvZHkiOnsibmFtZSI6ImJhcnJlZCJ9fX1dLFs0LDUsIiIsMSx7InN0eWxlIjp7ImJvZHkiOnsibmFtZSI6ImJhcnJlZCJ9fX1dLFs2LDddLFs3LDgsIlhee1xcb3B9IiwyLHsic3R5bGUiOnsiYm9keSI6eyJuYW1lIjoiYmFycmVkIn19fV0sWzksMTAsIiIsMCx7InN0eWxlIjp7ImJvZHkiOnsibmFtZSI6ImJhcnJlZCJ9fX1dLFsxMSw4XSxbNiw5LCJZXntcXG9wfSIsMCx7InN0eWxlIjp7ImJvZHkiOnsibmFtZSI6ImJhcnJlZCJ9fX1dLFsxMCwxMSwiWl57XFxvcH0iLDAseyJzdHlsZSI6eyJib2R5Ijp7Im5hbWUiOiJiYXJyZWQifX19XSxbMTIsMTcsImZcXGNvbG9uIFggXFxyaWdodGFycm93IFogXFxvdGltZXMgXFxjZG90cyBcXG90aW1lcyBZIiwxLHsic2hvcnRlbiI6eyJzb3VyY2UiOjIwLCJ0YXJnZXQiOjIwfX1dLFsyMCwxOSwiZl57XFxvcH0gXFxjb2xvbiBaXFxvdGltZXMgXFxjZG90cyBcXG90aW1lcyBZXFxyaWdodGFycm93IFgiLDEseyJzaG9ydGVuIjp7InNvdXJjZSI6MjAsInRhcmdldCI6MjB9fV0sWzE0LDE4LCIiLDAseyJzaG9ydGVuIjp7InNvdXJjZSI6MjAsInRhcmdldCI6MjB9LCJsZXZlbCI6MSwic3R5bGUiOnsidGFpbCI6eyJuYW1lIjoibWFwcyB0byJ9fX1dXQ==
		\[\begin{tikzcd}[ampersand replacement=\&,cramped,column sep=scriptsize]
			{C_1'} \& \cdots \& \cdots \& {C_2'} \&\& {(C_1')^{\op}} \& \cdots \& \cdots \& {(C_2')^{\op}} \\
			\\
			{C_1} \&\&\& {C_2} \&\& {C_1^{\op}} \&\&\& {C_2^{\op}}
			\arrow["Y"{inner sep=.8ex}, "\shortmid"{marking}, from=1-1, to=1-2]
			\arrow[""{name=0, anchor=center, inner sep=0}, "\shortmid"{marking}, from=1-2, to=1-3]
			\arrow["Z"{inner sep=.8ex}, "\shortmid"{marking}, from=1-3, to=1-4]
			\arrow["{Y^{\op}}"{inner sep=.8ex}, "\shortmid"{marking}, from=1-6, to=1-7]
			\arrow[""{name=1, anchor=center, inner sep=0}, from=1-6, to=3-6]
			\arrow[""{name=2, anchor=center, inner sep=0}, "\shortmid"{marking}, from=1-7, to=1-8]
			\arrow["{Z^{\op}}"{inner sep=.8ex}, "\shortmid"{marking}, from=1-8, to=1-9]
			\arrow[from=1-9, to=3-9]
			\arrow[from=3-1, to=1-1]
			\arrow[""{name=3, anchor=center, inner sep=0}, "X"'{inner sep=.8ex}, "\shortmid"{marking}, from=3-1, to=3-4]
			\arrow[""{name=4, anchor=center, inner sep=0}, from=3-4, to=1-4]
			\arrow[""{name=5, anchor=center, inner sep=0}, "{X^{\op}}"'{inner sep=.8ex}, "\shortmid"{marking}, from=3-6, to=3-9]
			\arrow["{f^{\op} \colon Z\otimes \cdots \otimes Y\rightarrow X}"{description}, between={0.2}{0.8}, Rightarrow, nfold, from=2, to=5]
			\arrow["{f\colon X \rightarrow Z \otimes \cdots \otimes Y}"{description}, between={0.2}{0.8}, Rightarrow, nfold, from=3, to=0]
			\arrow[between={0.2}{0.8}, maps to, from=4, to=1]
		\end{tikzcd}\]
	\end{description}
	The virtual double functor $\op^{t}_{\cV} \colon \op_{V}^{t}(\Mod(\cV)) \rightarrow \CoMod(\op_1(\cV))$ can be obtained similarly.
\end{proof}

Finally, we shall similarly extend the reversal duality. 

\begin{proposition}
	The $r$ morphism of involution on monoidal categories extends to the virtual double categories case through the following strict $2$-transformation cells:
	% https://q.uiver.app/#q=WzAsNCxbMCwwLCJcXE1vbkNhdF9sIl0sWzAsMiwiXFxNb25DYXRfbCJdLFsyLDAsIlxcdkRibENhdCJdLFsyLDIsIlxcdkRibENhdCJdLFswLDIsIlxcTW9kKC0pIl0sWzEsMywiXFxNb2QoLSkiLDJdLFswLDEsInIiLDJdLFsyLDMsIlxcb3BfSCJdLFsyLDEsInIiLDAseyJzaG9ydGVuIjp7InNvdXJjZSI6MjAsInRhcmdldCI6MjB9LCJsZXZlbCI6Mn1dXQ==
	\[\begin{tikzcd}[ampersand replacement=\&]
		{\MonCat_l} \&\& \vDblCat \\
		\\
		{\MonCat_l} \&\& \vDblCat
		\arrow["{\Mod(-)}", from=1-1, to=1-3]
		\arrow["r"', from=1-1, to=3-1]
		\arrow["r", between={0.2}{0.8}, Rightarrow, nfold, from=1-3, to=3-1]
		\arrow["{\op_H}", from=1-3, to=3-3]
		\arrow["{\Mod(-)}"', from=3-1, to=3-3]
	\end{tikzcd}\cdpunct[8pt]{.}\]
\end{proposition}

\begin{proof}
	As before, we specify the virtual double functor $r_{\cV}\colon \op_{H}(\Mod(\cV)) \rightarrow \Mod(r(\cV))$.
	\begin{description}
		\item[Objects] 
		% https://q.uiver.app/#q=WzAsMixbMCwwLCJNIFxcY29sb25lcXEgKE0sXFxtdSxlKSJdLFsyLDAsIk1eciBcXGNvbG9uZXFxIChNLCBcXG11LCBlKSJdLFswLDEsIiIsMCx7InN0eWxlIjp7InRhaWwiOnsibmFtZSI6Im1hcHMgdG8ifX19XV0=
		\[\begin{tikzcd}[ampersand replacement=\&]
			{M \coloneqq (M,\mu,e)} \&\& {M^r \coloneqq (M, \mu, e)}
			\arrow[maps to, from=1-1, to=1-3]
		\end{tikzcd}\]
		\item[V $1$-cell] 
		% https://q.uiver.app/#q=WzAsNCxbMCwwLCJNIl0sWzAsMiwiTSciXSxbMywwLCJNXnIiXSxbMywyLCIoTScpXnIiXSxbMCwxXSxbMiwzXSxbNCw1LCIiLDAseyJzaG9ydGVuIjp7InNvdXJjZSI6MjAsInRhcmdldCI6MjB9LCJsZXZlbCI6MSwic3R5bGUiOnsidGFpbCI6eyJuYW1lIjoibWFwcyB0byJ9fX1dXQ==
		\[\begin{tikzcd}[ampersand replacement=\&]
			M \&\&\& {M^r} \\
			\\
			{M'} \&\&\& {(M')^r}
			\arrow[""{name=0, anchor=center, inner sep=0}, from=1-1, to=3-1]
			\arrow[""{name=1, anchor=center, inner sep=0}, from=1-4, to=3-4]
			\arrow[between={0.2}{0.8}, maps to, from=0, to=1]
		\end{tikzcd}\]
		\item[H $1$-cell] 
		% https://q.uiver.app/#q=WzAsNCxbMCwwLCJNXzIiXSxbMywwLCJNXzEiXSxbNiwwLCIoTV8yKV5yIl0sWzksMCwiKE1fMSleciJdLFsxLDAsIlggXFxjb2xvbmVxcShYLFxcbGFtYmRhLFxccmhvKSIsMix7InN0eWxlIjp7ImJvZHkiOnsibmFtZSI6ImJhcnJlZCJ9fX1dLFsyLDMsIlhee3J9IFxcY29sb25lcXEgKFggLFxccmhvLFxcbGFtYmRhKSIsMCx7InN0eWxlIjp7ImJvZHkiOnsibmFtZSI6ImJhcnJlZCJ9fX1dLFs0LDUsIiIsMCx7InNob3J0ZW4iOnsic291cmNlIjo0MCwidGFyZ2V0Ijo0MH0sImxldmVsIjoxLCJzdHlsZSI6eyJ0YWlsIjp7Im5hbWUiOiJtYXBzIHRvIn19fV1d
		\[\begin{tikzcd}[ampersand replacement=\&]
			{M_2} \&\&\& {M_1} \&\&\& {(M_2)^r} \&\&\& {(M_1)^r}
			\arrow[""{name=0, anchor=center, inner sep=0}, "{X \coloneqq(X,\lambda,\rho)}"'{inner sep=.8ex}, "\shortmid"{marking}, from=1-4, to=1-1]
			\arrow[""{name=1, anchor=center, inner sep=0}, "{X^{r} \coloneqq (X ,\rho,\lambda)}"{inner sep=.8ex}, "\shortmid"{marking}, from=1-7, to=1-10]
			\arrow[between={0.4}{0.6}, maps to, from=0, to=1]
		\end{tikzcd}\]
		\item[Sq multicell] 
		% https://q.uiver.app/#q=WzAsMTIsWzAsMiwiTV8yIl0sWzMsMiwiTV8xIl0sWzAsMCwiTV8yJyJdLFszLDAsIk1fMSciXSxbMSwwLCJcXGNkb3RzIl0sWzIsMCwiXFxjZG90cyJdLFs1LDAsIihNXzInKV5yIl0sWzUsMiwiTV8yXnIiXSxbOCwyLCJNXzFeciJdLFs2LDAsIlxcY2RvdHMiXSxbNywwLCJcXGNkb3RzIl0sWzgsMCwiKE1fMScpXnIiXSxbMSwwLCJYIiwwLHsic3R5bGUiOnsiYm9keSI6eyJuYW1lIjoiYmFycmVkIn19fV0sWzIsMF0sWzMsMV0sWzQsMiwiWiIsMix7InN0eWxlIjp7ImJvZHkiOnsibmFtZSI6ImJhcnJlZCJ9fX1dLFszLDUsIlkiLDIseyJzdHlsZSI6eyJib2R5Ijp7Im5hbWUiOiJiYXJyZWQifX19XSxbNSw0LCIiLDEseyJzdHlsZSI6eyJib2R5Ijp7Im5hbWUiOiJiYXJyZWQifX19XSxbNiw3XSxbNyw4LCJYXnIiLDIseyJzdHlsZSI6eyJib2R5Ijp7Im5hbWUiOiJiYXJyZWQifX19XSxbOSwxMCwiIiwwLHsic3R5bGUiOnsiYm9keSI6eyJuYW1lIjoiYmFycmVkIn19fV0sWzExLDhdLFs2LDksIlpeciIsMCx7InN0eWxlIjp7ImJvZHkiOnsibmFtZSI6ImJhcnJlZCJ9fX1dLFsxMCwxMSwiWV5yIiwwLHsic3R5bGUiOnsiYm9keSI6eyJuYW1lIjoiYmFycmVkIn19fV0sWzIwLDE5LCJmIFxcY29sb24gWV5yIFxcb3RpbWVzX3tyfSBcXGNkb3RzIFxcb3RpbWVzX3tyfSBaXnIgXFxyaWdodGFycm93IFheciIsMSx7InNob3J0ZW4iOnsic291cmNlIjoyMCwidGFyZ2V0IjoyMH19XSxbMTQsMTgsIiIsMCx7InNob3J0ZW4iOnsic291cmNlIjoyMCwidGFyZ2V0IjoyMH0sImxldmVsIjoxLCJzdHlsZSI6eyJ0YWlsIjp7Im5hbWUiOiJtYXBzIHRvIn19fV0sWzE3LDEyLCJmXFxjb2xvbiBaIFxcb3RpbWVzIFxcY2RvdHMgXFxvdGltZXMgWSBcXHJpZ2h0YXJyb3cgWCIsMSx7InNob3J0ZW4iOnsic291cmNlIjoyMCwidGFyZ2V0IjoyMH19XV0=
		\[\begin{tikzcd}[ampersand replacement=\&,cramped]
			{M_2'} \& \cdots \& \cdots \& {M_1'} \&\& {(M_2')^r} \& \cdots \& \cdots \& {(M_1')^r} \\
			\\
			{M_2} \&\&\& {M_1} \&\& {M_2^r} \&\&\& {M_1^r}
			\arrow[from=1-1, to=3-1]
			\arrow["Z"'{inner sep=.8ex}, "\shortmid"{marking}, from=1-2, to=1-1]
			\arrow[""{name=0, anchor=center, inner sep=0}, "\shortmid"{marking}, from=1-3, to=1-2]
			\arrow["Y"'{inner sep=.8ex}, "\shortmid"{marking}, from=1-4, to=1-3]
			\arrow[""{name=1, anchor=center, inner sep=0}, from=1-4, to=3-4]
			\arrow["{Z^r}"{inner sep=.8ex}, "\shortmid"{marking}, from=1-6, to=1-7]
			\arrow[""{name=2, anchor=center, inner sep=0}, from=1-6, to=3-6]
			\arrow[""{name=3, anchor=center, inner sep=0}, "\shortmid"{marking}, from=1-7, to=1-8]
			\arrow["{Y^r}"{inner sep=.8ex}, "\shortmid"{marking}, from=1-8, to=1-9]
			\arrow[from=1-9, to=3-9]
			\arrow[""{name=4, anchor=center, inner sep=0}, "X"{inner sep=.8ex}, "\shortmid"{marking}, from=3-4, to=3-1]
			\arrow[""{name=5, anchor=center, inner sep=0}, "{X^r}"'{inner sep=.8ex}, "\shortmid"{marking}, from=3-6, to=3-9]
			\arrow["{f\colon Z \otimes \cdots \otimes Y \rightarrow X}"{description}, between={0.2}{0.8}, Rightarrow, nfold, from=0, to=4]
			\arrow[between={0.2}{0.8}, maps to, from=1, to=2]
			\arrow["{f \colon Y^r \otimes_{r} \cdots \otimes_{r} Z^r \rightarrow X^r}"{description}, between={0.2}{0.8}, Rightarrow, nfold, from=3, to=5]
		\end{tikzcd}\]
	\end{description}
\end{proof}

There are many other variants to consider, such as the pseudo case or the commutative case of each of these. However, the respective results are mostly analogous. 

\begin{comment}
	So, we shall only state that there is a relation between $\Mod()$ and $\cMod()$
	
	% https://q.uiver.app/#q=WzAsOCxbNiwwLCJcXE1vbkNhdF9sIl0sWzMsMCwiXFxCck1vbkNhdF9sIl0sWzAsMCwiXFxTeW1Nb25DYXRfbCJdLFszLDQsIlxcdkRibENhdCJdLFswLDUsIlxcU3ltTW9uQ2F0X2MiXSxbMyw1LCJcXEJyTW9uQ2F0X2MiXSxbNiw1LCJcXE1vbkNhdF9jIl0sWzMsOSwiXFx2RGJsQ2F0Il0sWzEsMCwiVSJdLFsyLDEsIlUiXSxbMCwzLCJcXE1vZCgtKSIsMV0sWzEsMywiXFxjTW9kKC0pIiwxXSxbMiwzLCJcXGNNb2QoLSkiLDFdLFs2LDcsIlxcQ29Nb2QoLSkiLDFdLFs1LDYsIlUiXSxbNCw1LCJVIl0sWzUsNywiXFxjQ29Nb2QoLSkiLDFdLFs0LDcsIlxcY0NvTW9kKC0pIiwxXSxbMTEsMTAsIiIsMSx7InNob3J0ZW4iOnsic291cmNlIjozMCwidGFyZ2V0IjozMH19XSxbMTIsMTEsIiIsMSx7InNob3J0ZW4iOnsic291cmNlIjozMCwidGFyZ2V0IjozMH0sInN0eWxlIjp7ImhlYWQiOnsibmFtZSI6Im5vbmUifX19XSxbMTcsMTYsIiIsMSx7InNob3J0ZW4iOnsic291cmNlIjozMCwidGFyZ2V0IjozMH0sInN0eWxlIjp7ImhlYWQiOnsibmFtZSI6Im5vbmUifX19XSxbMTYsMTMsIiIsMSx7InNob3J0ZW4iOnsic291cmNlIjozMCwidGFyZ2V0IjozMH19XV0=
	\[\begin{tikzcd}[ampersand replacement=\&]
		{\SymMonCat_l} \&\&\& {\BrMonCat_l} \&\&\& {\MonCat_l} \\
		\\
		\\
		\\
		\&\&\& \vDblCat \\
		{\SymMonCat_c} \&\&\& {\BrMonCat_c} \&\&\& {\MonCat_c} \\
		\\
		\\
		\\
		\&\&\& \vDblCat
		\arrow["U", from=1-1, to=1-4]
		\arrow[""{name=0, anchor=center, inner sep=0}, "{\cMod(-)}"{description}, from=1-1, to=5-4]
		\arrow["U", from=1-4, to=1-7]
		\arrow[""{name=1, anchor=center, inner sep=0}, "{\cMod(-)}"{description}, from=1-4, to=5-4]
		\arrow[""{name=2, anchor=center, inner sep=0}, "{\Mod(-)}"{description}, from=1-7, to=5-4]
		\arrow["U", from=6-1, to=6-4]
		\arrow[""{name=3, anchor=center, inner sep=0}, "{\cCoMod(-)}"{description}, from=6-1, to=10-4]
		\arrow["U", from=6-4, to=6-7]
		\arrow[""{name=4, anchor=center, inner sep=0}, "{\cCoMod(-)}"{description}, from=6-4, to=10-4]
		\arrow[""{name=5, anchor=center, inner sep=0}, "{\CoMod(-)}"{description}, from=6-7, to=10-4]
		\arrow[between={0.3}{0.7}, equals, from=0, to=1]
		\arrow[between={0.3}{0.7}, Rightarrow, from=1, to=2]
		\arrow[between={0.3}{0.7}, equals, from=3, to=4]
		\arrow[between={0.3}{0.7}, Rightarrow, from=4, to=5]
	\end{tikzcd}\cdpunct[8pt]{.}\]
\end{comment}

\begin{remark}
	In the case of a braided or symmetric monoidal category $\cV$, we can again compose with the canonical braiding functor $b_{\cV}$ to obtain the following inverse form of reversal component at $\cV$:
	% https://q.uiver.app/#q=WzAsMyxbMCwwLCJcXG9wX3tIfShcXE1vZChcXGNWKSkiXSxbMiwwLCJcXE1vZChyKFxcY1YpKSJdLFs0LDAsIlxcTW9kKFxcY1YpIl0sWzAsMSwicl97XFxjVn0iXSxbMSwyLCJcXE1vZChiX3tcXGNWfSkiXV0=
	\[\begin{tikzcd}[ampersand replacement=\&]
		{\op_{H}(\Mod(\cV))} \&\& {\Mod(r(\cV))} \&\& {\Mod(\cV)}
		\arrow["{r_{\cV}}", from=1-1, to=1-3]
		\arrow["{\Mod(b_{\cV})}", from=1-3, to=1-5]
	\end{tikzcd}\cdpunct[2pt]{.}\]
\end{remark}

Another well known construction on virtual double categories is the $2$-functor $\Mat \colon \MonCat_{l} \rightarrow \vDblCat$. Although this functor is not our central focus, it is worth discussing the dualities of this functor. We state the following lemma on reversal duality of $\Mat$.

\begin{lemma}
	The $2$-functor $\Mat$ extends to a $1$-morphism of involutions by the strict $2$-natural transformation
	% https://q.uiver.app/#q=WzAsNCxbMCwwLCJcXE1vbkNhdF9sIl0sWzAsMiwiXFxNb25DYXRfbCJdLFsyLDAsIlxcdkRibENhdCJdLFsyLDIsIlxcdkRibENhdCJdLFswLDIsIlxcTWF0KC0pIl0sWzEsMywiXFxNYXQoLSkiLDJdLFswLDEsInIiLDJdLFsyLDMsIlxcb3BfSCJdLFsyLDEsInIiLDAseyJzaG9ydGVuIjp7InNvdXJjZSI6MjAsInRhcmdldCI6MjB9LCJsZXZlbCI6Mn1dXQ==
	\[\begin{tikzcd}[ampersand replacement=\&,cramped]
		{\MonCat_l} \&\& \vDblCat \\
		\\
		{\MonCat_l} \&\& \vDblCat
		\arrow["{\Mat(-)}", from=1-1, to=1-3]
		\arrow["r"', from=1-1, to=3-1]
		\arrow["r", between={0.2}{0.8}, Rightarrow, nfold, from=1-3, to=3-1]
		\arrow["{\op_H}", from=1-3, to=3-3]
		\arrow["{\Mat(-)}"', from=3-1, to=3-3]
	\end{tikzcd}\cdpunct[8pt]{.}\]
\end{lemma}
\begin{proof}
	Let $(-)^t$ here denote the transpose of a matrix. The components $r_{\cV} \colon \op_H(\Mat(\cV)) \rightarrow \Mat(r(\cV))$ are defined by
	\begin{description}
		\item[Objects] 
		% https://q.uiver.app/#q=WzAsMixbMCwwLCJYIl0sWzIsMCwiWCJdLFswLDEsIiIsMCx7InN0eWxlIjp7InRhaWwiOnsibmFtZSI6Im1hcHMgdG8ifX19XV0=
		\[\begin{tikzcd}[ampersand replacement=\&,cramped]
			X \&\& X
			\arrow[maps to, from=1-1, to=1-3]
		\end{tikzcd}\]
		\item[V $1$-cell] 
		% https://q.uiver.app/#q=WzAsNCxbMCwwLCJYIl0sWzAsMiwiWSJdLFszLDAsIlgiXSxbMywyLCJZIl0sWzAsMV0sWzIsM10sWzQsNSwiIiwwLHsic2hvcnRlbiI6eyJzb3VyY2UiOjIwLCJ0YXJnZXQiOjIwfSwibGV2ZWwiOjEsInN0eWxlIjp7InRhaWwiOnsibmFtZSI6Im1hcHMgdG8ifX19XV0=
		\[\begin{tikzcd}[ampersand replacement=\&,cramped]
			X \&\&\& X \\
			\\
			Y \&\&\& Y
			\arrow[""{name=0, anchor=center, inner sep=0}, from=1-1, to=3-1]
			\arrow[""{name=1, anchor=center, inner sep=0}, from=1-4, to=3-4]
			\arrow[between={0.2}{0.8}, maps to, from=0, to=1]
		\end{tikzcd}\]
		\item[H $1$-cell] 
		% https://q.uiver.app/#q=WzAsNCxbMCwwLCJZIl0sWzMsMCwiWCJdLFs2LDAsIlkiXSxbOSwwLCJYIl0sWzEsMCwiTSIsMix7InN0eWxlIjp7ImJvZHkiOnsibmFtZSI6ImJhcnJlZCJ9fX1dLFsyLDMsIk1edCIsMCx7InN0eWxlIjp7ImJvZHkiOnsibmFtZSI6ImJhcnJlZCJ9fX1dLFs0LDUsIiIsMCx7InNob3J0ZW4iOnsic291cmNlIjo0MCwidGFyZ2V0Ijo0MH0sImxldmVsIjoxLCJzdHlsZSI6eyJ0YWlsIjp7Im5hbWUiOiJtYXBzIHRvIn19fV1d
		\[\begin{tikzcd}[ampersand replacement=\&,cramped]
			Y \&\&\& X \&\&\& Y \&\&\& X
			\arrow[""{name=0, anchor=center, inner sep=0}, "M"'{inner sep=.8ex}, "\shortmid"{marking}, from=1-4, to=1-1]
			\arrow[""{name=1, anchor=center, inner sep=0}, "{M^t}"{inner sep=.8ex}, "\shortmid"{marking}, from=1-7, to=1-10]
			\arrow[between={0.4}{0.6}, maps to, from=0, to=1]
		\end{tikzcd}\]
		\item[Sq multicell] 
		% https://q.uiver.app/#q=WzAsMTIsWzAsMiwiWV8xIl0sWzMsMiwiWV8wIl0sWzAsMCwiWF9uIl0sWzMsMCwiWF8wIl0sWzEsMCwiXFxjZG90cyJdLFsyLDAsIlxcY2RvdHMiXSxbNSwwLCJYX24iXSxbNSwyLCJZXzEiXSxbOCwyLCJNXzFeciJdLFs2LDAsIlxcY2RvdHMiXSxbNywwLCJcXGNkb3RzIl0sWzgsMCwiWF8wIl0sWzEsMCwiTiIsMCx7InN0eWxlIjp7ImJvZHkiOnsibmFtZSI6ImJhcnJlZCJ9fX1dLFsyLDBdLFszLDFdLFs0LDIsIk1fbiIsMix7InN0eWxlIjp7ImJvZHkiOnsibmFtZSI6ImJhcnJlZCJ9fX1dLFszLDUsIk1fMSIsMix7InN0eWxlIjp7ImJvZHkiOnsibmFtZSI6ImJhcnJlZCJ9fX1dLFs1LDQsIiIsMSx7InN0eWxlIjp7ImJvZHkiOnsibmFtZSI6ImJhcnJlZCJ9fX1dLFs2LDddLFs3LDgsIk5edCIsMix7InN0eWxlIjp7ImJvZHkiOnsibmFtZSI6ImJhcnJlZCJ9fX1dLFs5LDEwLCIiLDAseyJzdHlsZSI6eyJib2R5Ijp7Im5hbWUiOiJiYXJyZWQifX19XSxbMTEsOF0sWzYsOSwiTV9uXnQiLDAseyJzdHlsZSI6eyJib2R5Ijp7Im5hbWUiOiJiYXJyZWQifX19XSxbMTAsMTEsIk1fMV50IiwwLHsic3R5bGUiOnsiYm9keSI6eyJuYW1lIjoiYmFycmVkIn19fV0sWzIwLDE5LCIiLDEseyJzaG9ydGVuIjp7InNvdXJjZSI6MjAsInRhcmdldCI6MjB9fV0sWzE0LDE4LCIiLDAseyJzaG9ydGVuIjp7InNvdXJjZSI6MjAsInRhcmdldCI6MjB9LCJsZXZlbCI6MSwic3R5bGUiOnsidGFpbCI6eyJuYW1lIjoibWFwcyB0byJ9fX1dLFsxNywxMiwiIiwxLHsic2hvcnRlbiI6eyJzb3VyY2UiOjIwLCJ0YXJnZXQiOjIwfX1dXQ==
		\[\begin{tikzcd}[ampersand replacement=\&,cramped]
			{X_n} \& \cdots \& \cdots \& {X_0} \&\& {X_n} \& \cdots \& \cdots \& {X_0} \\
			\\
			{Y_1} \&\&\& {Y_0} \&\& {Y_1} \&\&\& {M_1^r}
			\arrow[from=1-1, to=3-1]
			\arrow["{M_n}"'{inner sep=.8ex}, "\shortmid"{marking}, from=1-2, to=1-1]
			\arrow[""{name=0, anchor=center, inner sep=0}, "\shortmid"{marking}, from=1-3, to=1-2]
			\arrow["{M_1}"'{inner sep=.8ex}, "\shortmid"{marking}, from=1-4, to=1-3]
			\arrow[""{name=1, anchor=center, inner sep=0}, from=1-4, to=3-4]
			\arrow["{M_n^t}"{inner sep=.8ex}, "\shortmid"{marking}, from=1-6, to=1-7]
			\arrow[""{name=2, anchor=center, inner sep=0}, from=1-6, to=3-6]
			\arrow[""{name=3, anchor=center, inner sep=0}, "\shortmid"{marking}, from=1-7, to=1-8]
			\arrow["{M_1^t}"{inner sep=.8ex}, "\shortmid"{marking}, from=1-8, to=1-9]
			\arrow[from=1-9, to=3-9]
			\arrow[""{name=4, anchor=center, inner sep=0}, "N"{inner sep=.8ex}, "\shortmid"{marking}, from=3-4, to=3-1]
			\arrow[""{name=5, anchor=center, inner sep=0}, "{N^t}"'{inner sep=.8ex}, "\shortmid"{marking}, from=3-6, to=3-9]
			\arrow[between={0.2}{0.8}, Rightarrow, nfold, from=0, to=4]
			\arrow[between={0.2}{0.8}, maps to, from=1, to=2]
			\arrow[between={0.2}{0.8}, Rightarrow, nfold, from=3, to=5]
		\end{tikzcd}\]
	\end{description}
\end{proof}

\begin{remark}
	More generally, \cite{fujii2025famenr} propose a generalized matrices functor $\Mat \colon \vDblCat \rightarrow \vDblCat$, which arises primarily from the discrete opfibration $\Set_* \rightarrow \Set$. In particular, they define it as a composite functor % https://q.uiver.app/#q=WzAsNCxbMCwwLCJcXHZEYmxDYXQiXSxbMywwLCJcXHZEYmxDYXQvKFxcU2V0XyopX3tcXG1hdGhybXtoY319Il0sWzMsMiwiXFx2RGJsQ2F0LyhcXFNldClfe1xcbWF0aHJte2hjfX0iXSxbMCwyLCJcXHZEYmxDYXQiXSxbMSwyLCJcXHByb2Rfe1Bfe1xcbWF0aHJte2hjfX19Il0sWzIsMywiVSJdLFswLDEsIihcXFNldF8qKV97XFxtYXRocm17aGN9fSBcXHRpbWVzICgtKSJdLFswLDMsIlxcTWF0IiwyXV0=
	\[\begin{tikzcd}[ampersand replacement=\&,cramped]
		\vDblCat \&\&\& {\vDblCat/(\Set_*)_{\mathrm{hc}}} \\
		\\
		\vDblCat \&\&\& {\vDblCat/(\Set)_{\mathrm{hc}}}
		\arrow["{(\Set_*)_{\mathrm{hc}} \times (-)}", from=1-1, to=1-4]
		\arrow["\Mat"', from=1-1, to=3-1]
		\arrow["{\prod_{P_{\mathrm{hc}}}}", from=1-4, to=3-4]
		\arrow["U", from=3-4, to=3-1]
	\end{tikzcd}\cdpunct[8pt]{.}\]
	One could show that each of the individual $2$-functor appearing in this decomposition is compatible with the $\op_{H}$ duality. In particular, this would generalize the previous lemma to $\vDblCat$. 
\end{remark}

Lastly, in contrast to the horizontal duality, we note that the $2$-functor $\Mat$ doesn't respect the vertical duality of virtual double categories in the expected way. Although one could define a suitable $2$-functor $\mathbf{CoMat}\colon \vcoDblCat \rightarrow \vcoDblCat$, applying $\CoMod$ to this construction would not yield the desired variance for morphisms between cocategories. This is more readily observed by considering the category $\cV^{\op}\Cat$ of $\op_1(\cV)$-enriched categories. In particular, the action of functors on the hom-objects and on the set of objects are of opposite variance. If one is content with this disparity, we believe that one could dualize the $\Mat$ construction.

\subsection{Higher Dualities} \label{sub:highermonduals}

We briefly describe the analogous transport phenomenon in the case of $\nCat{1}$-enriched symmetric monoidal category $\nCat{1}$. 

\begin{definition}
	Let $(\cV,\otimes, I)$ be any monoidal category. We define a left $\cV$-module category $\cC$ as a left pseudomodule object in $\nCat{1}$. The definitions for right module or bimodule categories are analogous.
\end{definition}

To describe the $\op_1$ duality of monoidal categories, we need to have the variants $\nMon{1}_c$ and $\nMon{1}_l$ corresponding the construction $\nMon{1}$. This is similar to how one dimension lower, we required separate $\Mon$ and $\CoMon$ constructions. In particular, $\nMon{1}_l$ is simply the construction we have already provided involving lax $1$-morphisms. Whereas, $\nMon{1}_c$ is the colax $1$-morphism variant. 
 
\begin{lemma}
	The $\op_1$ and $\op_1^t$ duality inverses of monoidal categories are equivalently expressed as the two maps:
	% https://q.uiver.app/#q=WzAsOCxbMSwwLCJcXG9wXzIoXFxuTW9uezF9X3tsfShcXG5DYXR7MX0pKSJdLFszLDAsIlxcbk1vbnsxfV97Y30oXFxvcF8yKFxcbkNhdHsxfSkpIl0sWzEsMSwiXFxvcF97Mn0oXFxuTW9uezF9X3tjfShcXG5DYXR7MX0pKSJdLFszLDEsIlxcbk1vbnsxfV97bH0oXFxvcF8yKFxcbkNhdHsxfSkpIl0sWzYsMCwiXFxuTW9uezF9X3tjfShcXG5DYXR7MX0pIl0sWzYsMSwiXFxuTW9uezF9X3tsfShcXG5DYXR7MX0pIl0sWzAsMCwiXFxvcF8xXnt0fSBcXGNvbG9uIl0sWzAsMSwiXFxvcF8xIFxcY29sb24iXSxbMCwxXSxbMiwzXSxbMSw0LCJcXG5Nb257MX1fe2N9KFxcb3BfMSkiXSxbMyw1LCJcXG5Nb257MX1fe2x9KFxcb3BfMSkiXV0=
	\[\begin{tikzcd}[ampersand replacement=\&,column sep=scriptsize]
		{\op_1^{t} \colon} \& {\op_2(\nMon{1}_{l}(\nCat{1}))} \&\& {\nMon{1}_{c}(\op_2(\nCat{1}))} \&\&\& {\nMon{1}_{c}(\nCat{1})} \\
		{\op_1 \colon} \& {\op_{2}(\nMon{1}_{c}(\nCat{1}))} \&\& {\nMon{1}_{l}(\op_2(\nCat{1}))} \&\&\& {\nMon{1}_{l}(\nCat{1})}
		\arrow[from=1-2, to=1-4]
		\arrow["{\nMon{1}_{c}(\op_1)}", from=1-4, to=1-7]
		\arrow[from=2-2, to=2-4]
		\arrow["{\nMon{1}_{l}(\op_1)}", from=2-4, to=2-7]
	\end{tikzcd}\]
\end{lemma}
 
\begin{proof}
	The identifications $\nMon{1}_{l}(\op_2(\nCat{1})) \cong \op_2(\MonCat_{c})$ and $\nMon{1}_{c}(\op_2(\nCat{1})) \cong \op_2(\MonCat_{l})$ follows directly by the definitions. One should view these identifications as merely comprehension identifications. Whereas, $\op_1 : \op_2(\nCat{1}) \rightarrow \nCat{1}$ is a strict monoidal $2$-functor. Thus, we have transport morphisms $\nMon{1}_c(\op_1)$ and $\nMon{1}_l(\op_1)$ given by applying the $\op_1$ functor directly on data of the objects and the morphisms. This is exactly the description we provided at the beginning of this section. Now, we are stating the same description in a more formalized way.
\end{proof}

In other words, we have fully expressed the $\op_1$ and $\op_1^t$ duality of monoidal categories in terms of a hypothetical higher duality $\op_2\colon \op_3(\mathbf{MonBicat_{x}}) \rightarrow \mathbf{MonBicat_{y}}$, where $x$ and $y$ denote some suitable variants. More precisely, assuming we had all the relevant definitions, the duality can be described by the following diagram
% https://q.uiver.app/#q=WzAsNixbMiwwLCJcXG9wXzMoXFxtYXRoYmZ7TW9uQmljYXR9X3gpIl0sWzIsMiwiXFxtYXRoYmZ7TW9uQmljYXR9X3kiXSxbNCwwLCJcXG9wXzMoXFxCaUNhdF94KSJdLFs0LDIsIlxcQmlDYXRfeSJdLFswLDAsIioiXSxbMCwyLCIqIl0sWzAsMSwiXFxvcF8yIiwxXSxbMCwyLCJcXG5Nb257MX1fbCJdLFsyLDMsIlxcb3BfMiJdLFsxLDMsIlxcbk1vbnsxfV9sIiwyXSxbMiwxLCJcXG9wXzIiLDEseyJzaG9ydGVuIjp7InNvdXJjZSI6MjAsInRhcmdldCI6MjB9LCJsZXZlbCI6Mn1dLFs0LDVdLFs1LDEsIlxcQ2F0IiwyXSxbNCwwLCJcXENhdCJdLFswLDUsIlxcb3BfMSIsMSx7InNob3J0ZW4iOnsic291cmNlIjoyMCwidGFyZ2V0IjoyMH0sImxldmVsIjoyfV1d
\[\begin{tikzcd}[ampersand replacement=\&]
	{*} \&\& {\op_3(\mathbf{MonBicat}_x)} \&\& {\op_3(\BiCat_x)} \\
	\\
	{*} \&\& {\mathbf{MonBicat}_y} \&\& {\BiCat_y}
	\arrow["\Cat", from=1-1, to=1-3]
	\arrow[from=1-1, to=3-1]
	\arrow["{\nMon{1}_l}", from=1-3, to=1-5]
	\arrow["{\op_1}"{description}, between={0.2}{0.8}, Rightarrow, nfold, from=1-3, to=3-1]
	\arrow["{\op_2}"{description}, from=1-3, to=3-3]
	\arrow["{\op_2}"{description}, between={0.2}{0.8}, Rightarrow, nfold, from=1-5, to=3-3]
	\arrow["{\op_2}", from=1-5, to=3-5]
	\arrow["\Cat"', from=3-1, to=3-3]
	\arrow["{\nMon{1}_l}"', from=3-3, to=3-5]
\end{tikzcd}\cdpunct[8pt]{.}\]

It is now also straightforward to guess how pseudomodules in $\nCat{1}$ ought to transport under the dualities.

\begin{lemma}
	Let $\cV_1$ and $\cV_2$ be any two monoidal categories and let $\mathbf{Bimod}(\cV_1,\cV_2)$ be $2$-category of modules between $\cV_1$ and $\cV_2$. There exists the following involutive isomorphism induced by the action of $\op_1 \colon \op_2(\nCat{1}) \rightarrow \nCat{1}$:
	% https://q.uiver.app/#q=WzAsMyxbMSwwLCJcXG9wXzIoXFxtYXRoYmZ7Qmltb2R9KFxcY1ZfMSxcXGNWXzIpKSJdLFszLDAsIlxcbWF0aGJme0JpbW9kfShcXG9wXzEoXFxjVl8xKSxcXG9wXzEoXFxjVl8yKSkiXSxbMCwwLCJcXG9wXzEgXFxjb2xvbiJdLFswLDFdXQ==
	\[\begin{tikzcd}[ampersand replacement=\&,cramped]
		{\op_1 \colon} \& {\op_2(\mathbf{Bimod}(\cV_1,\cV_2))} \&\& {\mathbf{Bimod}(\op_1(\cV_1),\op_1(\cV_2))}
		\arrow[from=1-2, to=1-4]
	\end{tikzcd}\cdpunct[0pt]{.}\]
\end{lemma}

Simply note that all the defining axioms are comprised of natural isomorphisms which can be inverted. Finally, the duality for reversal is expressible similarly.

\begin{lemma} \leavevmode
	\begin{enumerate}
		\item The reversal duality $r : \MonCat_{x} \rightarrow \MonCat_{x}$ is given by the map
		% https://q.uiver.app/#q=WzAsNCxbMCwwLCJyIFxcY29sb24gIl0sWzEsMCwiXFxuTW9uezF9KFxcbkNhdHsxfSkiXSxbMywwLCJcXG5Nb24xKHIoXFxuQ2F0ezF9KSkiXSxbNiwwLCJcXG5Nb257MX0oXFxuQ2F0ezF9KSJdLFsxLDJdLFsyLDMsIlxcbk1vbnsxfSgoMSxzX3tcXG5DYXR7MX19LDEpKSJdXQ==
		\[\begin{tikzcd}[ampersand replacement=\&,cramped]
			{r \colon } \& {\nMon{1}(\nCat{1})} \&\& {\nMon1(r(\nCat{1}))} \&\&\& {\nMon{1}(\nCat{1})}
			\arrow[from=1-2, to=1-4]
			\arrow["{\nMon{1}((1,s_{\nCat{1}},1))}", from=1-4, to=1-7]
		\end{tikzcd}\cdpunct[0pt]{,}\]
		where $(1,s_{\nCat{1}},1)$ is the canonical identity symmetric monoidal strict $2$-functor whose compositor constraint is the symmetry of $\nCat{1}$. 
		\item There exists the following involutive isomorphism on the bimodule $2$-categories induced by the action of $r$:
		% https://q.uiver.app/#q=WzAsMyxbMCwwLCJyXFxjb2xvbiJdLFsxLDAsIlxcbWF0aGJme0JpbW9kfShcXGNWXzEsXFxjVl8yKSJdLFszLDAsIlxcbWF0aGJme0JpbW9kfShyKFxcY1ZfMikscihcXGNWXzEpKSJdLFsxLDJdXQ==
		\[\begin{tikzcd}[ampersand replacement=\&,cramped]
			{r\colon} \& {\mathbf{Bimod}(\cV_1,\cV_2)} \&\& {\mathbf{Bimod}(r(\cV_2),r(\cV_1))}
			\arrow[from=1-2, to=1-4]
		\end{tikzcd}\cdpunct[0pt]{.}\]
	\end{enumerate}
\end{lemma}

At this point, we could also obtain similar results relating monoidal and comonoidal categories using the $\op_1$ involution of bicategories, that is, $\op_1(\mathbf{1CoMon}_l(\nCat{1})) \cong \nMon{1}_l(\op_1(\nCat{1}))$. However, these results are not essential for us, so we do not investigate them here. 

